\providecommand{\N}{\mathbb{N}}
\providecommand{\tr}{\operatorname{tr}}
\providecommand{\Lap}{\mathcal{L}}
\providecommand{\Nbr}{\mathcal{N}}
\providecommand{\HASH}{\operatorname{HASH}}
\providecommand{\WL}{\textsc{wl}}
\providecommand{\Hess}{\operatorname{Hess}}
\providecommand{\dive}{\operatorname{div}}
\providecommand{\grad}{\operatorname{grad}}
\providecommand{\im}{\operatorname{im}}
\providecommand{\rank}{\operatorname{rank}}

\chapter{Graphs: Graph Neural Networks}
\label{ch:graphs}

Grids and groups, the subjects of \Cref{ch:grids,ch:groups}, share a common
luxury: a rigid, homogeneous domain in which every point looks like every other
point and the symmetry group acts transitively. A pixel grid is the same
everywhere; a sphere looks the same after any rotation. Most relational data of
practical interest enjoys no such regularity. A molecule, a social network, a
citation graph, a road map, or a triangulated 3D surface is an irregular
collection of entities tied together by pairwise relations, in which different
nodes have different numbers of neighbours and there is no canonical ordering of
those neighbours. The natural domain for such data is a \emph{graph}.

Graphs occupy a privileged position in geometric deep learning because they are
the bridge between the discrete world in which neural networks operate and the
continuous geometries (curved surfaces, three-dimensional Euclidean space, the
terrestrial sphere) in which much real data lives. Through \emph{proximity
graphs}, a continuous manifold is approximated by a finite set of points joined
by edges that encode neighbourhood relations, preserving essential topological
and spectral properties. This has two consequences. On the one hand it justifies,
theoretically, the use of graph neural networks for data that naturally lives on
manifolds (molecules, 3D point clouds, climate data). On the other it sets a
fundamental limit: any geometric information of the continuous structure that the
discretization fails to capture is inaccessible to graph-based architectures.

This chapter develops the third of the five Gs. We begin by viewing graphs as
\emph{discretizations of geometry} --- weighted graphs, $\epsilon$-neighbourhood
graphs, grids and triangulated surfaces as special cases, and the simplicial
machinery that underlies them. We then study the \emph{topological geometry} of
graphs --- connectivity, cycles, trees, and above all the graph Laplacian, the
discrete counterpart of the Laplace--Beltrami operator (\Cref{ch:geodesics}). A
section on \emph{topological deep learning} lifts message passing from nodes to
simplices through the boundary operator and the Hodge Laplacian. We make precise
the sense in which a graph \emph{converges} to a manifold, and we develop the
\emph{spectral analysis} of graphs --- the graph Fourier transform and spectral
filters. We then turn to the symmetry that governs graphs ---
\emph{permutation invariance} --- and the local computation it forces:
\emph{message passing}, the single mechanism that unifies the architectures of
the previous chapters. A convolutional network is message passing on a grid,
self-attention is message passing on a complete graph, and a graph neural
network is message passing on an arbitrary graph; we prove the first
identification with a worked numerical example. Finally we confront the question
of how \emph{expressive} these models are, exposing the fundamental limits set by
the Weisfeiler--Leman hierarchy and the companion pathologies of oversmoothing
and oversquashing.

\section{Graphs as discretizations of manifolds}
\label{sec:graphs:discretization}

A graph $G = (V, E)$ consists of a finite set of vertices (nodes)
$V = \{1, \dots, n\}$ and a set of edges $E \subseteq V \times V$. We attach to
each node a feature vector and collect them into a matrix
$X \in \R^{n \times d}$ whose $i$-th row $x_i \in \R^d$ describes node $i$. To
encode metric or relational information we equip the edges with weights.

\begin{definition}[Weighted graph]
\label{def:graphs:weighted}
A \emph{weighted graph} is a graph $G = (V, E)$ equipped with a weight function
$w \colon E \to \R$ assigning a numerical value to each edge. When the graph is
undirected we require $w(u,v) = w(v,u)$.
\end{definition}

\begin{definition}[Neighbourhood]
\label{def:graphs:neighbourhood}
For a vertex $v \in V$, the \emph{neighbourhood} of $v$ is
$\Nbr(v) = \{u \in V : (v,u) \in E\}$, the set of all vertices directly connected
to $v$ by an edge. Its \emph{degree} is $\deg(v) = \abs{\Nbr(v)}$ (or, for
weighted graphs, $\deg(v) = \sum_{u} w(v,u)$).
\end{definition}

The discrete approximation of a manifold $\mathcal{M}$ is built from a finite
sample of points together with a connectivity rule based on distances.

\begin{definition}[Neighbourhood graph]
\label{def:graphs:epsilon}
Given a set of points $\{x_i\}_{i=1}^n$ on a manifold $\mathcal{M}$, the
\emph{$\epsilon$-neighbourhood graph} $G = (V, E)$ has
\begin{itemize}
  \item $V = \{x_1, \dots, x_n\}$ (the vertices are the sampled points),
  \item $(x_i, x_j) \in E$ whenever $d_{\mathcal{M}}(x_i, x_j) < \epsilon$ for a
        threshold $\epsilon > 0$,
\end{itemize}
where $d_{\mathcal{M}}$ denotes the geodesic distance on $\mathcal{M}$
(\Cref{ch:geodesics}).
\end{definition}

This simple construction has profound consequences: if the points $\{x_i\}$ are
sufficiently dense in $\mathcal{M}$ and the radius $\epsilon$ is chosen
appropriately, the resulting graph $G$ captures fundamental topological and
geometric properties of the original manifold. \Cref{fig:graphs:epsilon}
illustrates the construction.

\begin{figure}[h]
\centering
\begin{tikzpicture}[font=\small,
    pnode/.style={circle, fill=gdlblue!70, inner sep=0pt, minimum size=4pt}]
% --- left: sampled points on a manifold ---
\begin{scope}[shift={(0,0)}]
  \begin{scope}[on background layer]
    \fill[gdlorange!12, rounded corners=3pt] (-1.6,-1.6) rectangle (1.6,1.6);
  \end{scope}
  \foreach \a/\r in {15/1.2, 55/1.35, 95/1.15, 140/1.3, 185/1.25, 225/1.1,
                     265/1.35, 305/1.2, 345/1.3, 35/0.6, 130/0.7, 250/0.55,
                     320/0.65, 200/0.4}
    \fill[gdlblue!70] (\a:\r) circle (2.2pt);
  \node[font=\bfseries, text=gdlorange!70!black] at (0,-2.0) {Sampled points};
\end{scope}
% --- right: epsilon graph ---
\begin{scope}[shift={(5,0)}]
  \begin{scope}[on background layer]
    \fill[gdlgreen!10, rounded corners=3pt] (-1.6,-1.6) rectangle (1.6,1.6);
  \end{scope}
  \foreach \i/\a/\r in {1/15/1.2, 2/55/1.35, 3/95/1.15, 4/140/1.3, 5/185/1.25,
                        6/225/1.1, 7/265/1.35, 8/305/1.2, 9/345/1.3, 10/35/0.6,
                        11/130/0.7, 12/250/0.55, 13/320/0.65, 14/200/0.4}
    \coordinate (p\i) at (\a:\r);
  \foreach \a/\b in {1/2,2/3,3/4,4/5,5/6,6/7,7/8,8/9,9/1,1/10,2/10,3/11,4/11,
                     6/12,7/12,8/13,9/13,5/14,6/14}
    \draw[gdlgreen!55, thick] (p\a) -- (p\b);
  \foreach \i in {1,...,14} \fill[gdlblue!70] (p\i) circle (2.2pt);
  \node[font=\bfseries, text=gdlgreen!70!black] at (0,-2.0) {$\epsilon$-graph};
\end{scope}
\draw[->, line width=1.4pt, gdlgray!70] (1.9,0) -- (3.1,0)
  node[midway, above, font=\scriptsize]{$d_{\mathcal{M}}<\epsilon$};
\end{tikzpicture}
\caption[$\epsilon$-neighbourhood graph]{Construction of an
$\epsilon$-neighbourhood graph. Points sampled from a manifold (left) are
connected whenever their distance is below $\epsilon$, forming a graph that
approximates the local structure of the original manifold (right).}
\label{fig:graphs:epsilon}
\end{figure}

\paragraph{Regular grids as a special case.}
The regular grids on which traditional convolutional networks operate are
particular instances of this general construction.

\begin{example}[2D image as a regular graph]
\label{ex:graphs:image-grid}
A digital image of size $H \times W$ may be regarded as a graph $G = (V, E)$ with
\begin{itemize}
  \item $V = \{(i,j) : 1 \le i \le H, \ 1 \le j \le W\}$ (pixels),
  \item $(i,j) \sim (i',j')$ whenever $\abs{i-i'} + \abs{j-j'} = 1$
        (4-neighbour connectivity).
\end{itemize}
This graph is a \emph{uniform discretization} of the continuous domain
$[0,1]\times[0,1]$ with the Euclidean metric. The regularity of the grid
translates into special properties: every interior node has degree exactly $4$;
the graph possesses a discrete translation symmetry; and the graph geodesic
distance approximates the $L^1$ distance in the continuous domain.
\end{example}

\begin{example}[Triangulated 3D surface as an irregular graph]
\label{ex:graphs:triangulated}
Consider a three-dimensional surface $\mathcal{S} \subset \R^3$ (for instance, a
3D model of an object). A \emph{triangulation} of $\mathcal{S}$ produces a graph
$G = (V, E)$ whose vertices $V$ are the triangulation vertices and whose edges
$E$ are the triangle edges. Unlike regular grids: node degrees are
\emph{heterogeneous} (typically between $4$ and $8$ for well-formed meshes);
there is no global translation symmetry; and the local curvature of $\mathcal{S}$
is reflected in the vertex density and graph irregularity. This construction is
fundamental in computer graphics, physical simulation and 3D shape analysis.
Since $\mathcal{S}$ is a surface (a compact 2-manifold), the existence of such a
triangulation is guaranteed by Radó's theorem (\Cref{thm:graphs:rado}). Graph
neural networks on these triangulations must respect \emph{gauge equivariance}
(\Cref{ch:gauges}): predictions cannot depend on the specific orientation of the
mesh or on the chosen parametrization.
\end{example}

\paragraph{Simplices, triangulations and simplicial complexes.}
To formalize triangulations and their generalizations we need the fundamental
building blocks of combinatorial topology.

\begin{definition}[$k$-simplex]
\label{def:graphs:simplex}
A \emph{$k$-simplex} $\sigma$ is the convex hull (the set of all convex
combinations) of $k+1$ affinely independent points $v_0, v_1, \dots, v_k$ in
$\R^n$:
\[
  \sigma = [v_0, v_1, \dots, v_k]
    = \Big\{ \textstyle\sum_{i=0}^k t_i v_i : \sum_{i=0}^k t_i = 1,\ t_i \ge 0 \Big\}.
\]
The points $v_i$ are its \emph{vertices} and $k$ is its \emph{dimension}. A
$0$-simplex is a point, a $1$-simplex an edge, a $2$-simplex a triangle, a
$3$-simplex a tetrahedron, and so on.
\end{definition}

\begin{definition}[Face of a simplex]
\label{def:graphs:face}
A \emph{face} of dimension $j < k$ of the $k$-simplex
$\sigma = [v_0, \dots, v_k]$ is any simplex spanned by a subset of $j+1$
vertices, $\tau = [v_{i_0}, \dots, v_{i_j}]$ with
$\{i_0, \dots, i_j\} \subseteq \{0, \dots, k\}$. We write $\tau \le \sigma$ when
$\tau$ is a face of $\sigma$.
\end{definition}

\begin{definition}[Simplicial complex]
\label{def:graphs:complex}
A \emph{simplicial complex} $\mathcal{K}$ is a finite collection of simplices
such that
\begin{enumerate}
  \item if $\sigma \in \mathcal{K}$ and $\tau \le \sigma$, then
        $\tau \in \mathcal{K}$ (closed under faces),
  \item if $\sigma, \tau \in \mathcal{K}$, then $\sigma \cap \tau$ is a face of
        both (well-defined intersection).
\end{enumerate}
The \emph{$1$-skeleton} of $\mathcal{K}$ is the graph $G = (V, E)$ formed by the
$0$-simplices (vertices) and $1$-simplices (edges).
\end{definition}

\begin{figure}[h]
\centering
\begin{tikzpicture}[>=Stealth, font=\small]
% --- top row: individual simplices ---
\node[font=\bfseries] at (-5, 3.8) {$0$-simplex};
\node[font=\bfseries] at (-1.5, 3.8) {$1$-simplex};
\node[font=\bfseries] at (2, 3.8) {$2$-simplex};
\node[font=\bfseries] at (6, 3.8) {$3$-simplex};
\fill[gdlblue!70] (-5, 2.5) circle (3pt);
\node[below, font=\scriptsize] at (-5, 2.3) {$v_0$};
\fill[gdlblue!70] (-2.3, 2.5) circle (3pt);
\fill[gdlblue!70] (-0.7, 2.5) circle (3pt);
\draw[thick, gdlblue!60] (-2.3, 2.5) -- (-0.7, 2.5);
\coordinate (t0) at (1, 2.2); \coordinate (t1) at (3, 2.2);
\coordinate (t2) at (2, 3.4);
\fill[gdlorange!30] (t0) -- (t1) -- (t2) -- cycle;
\draw[thick, gdlorange!70] (t0) -- (t1) -- (t2) -- cycle;
\foreach \p in {t0,t1,t2} \fill[gdlblue!70] (\p) circle (3pt);
\coordinate (s0) at (5, 2.1); \coordinate (s1) at (7, 2.1);
\coordinate (s2) at (6.4, 3.5); \coordinate (s3) at (5.8, 2.9);
\fill[gdlred!15, opacity=0.5] (s0) -- (s1) -- (s3) -- cycle;
\fill[gdlred!18, opacity=0.5] (s0) -- (s2) -- (s3) -- cycle;
\fill[gdlred!25, opacity=0.6] (s0) -- (s1) -- (s2) -- cycle;
\fill[gdlred!20, opacity=0.6] (s1) -- (s2) -- (s3) -- cycle;
\draw[thick, gdlred!50] (s0) -- (s1) -- (s2) -- cycle;
\draw[thick, gdlred!50] (s0) -- (s3); \draw[thick, gdlred!50] (s1) -- (s3);
\draw[thick, gdlred!50] (s2) -- (s3);
\foreach \p in {s0,s1,s2,s3} \fill[gdlblue!70] (\p) circle (3pt);
% --- separator ---
\draw[dashed, gdlgray!50] (-6.5, 1.2) -- (8, 1.2);
\node[font=\bfseries] at (0.5, 0.7) {Simplicial complex $\mathcal{K}$};
% --- bottom row: assembled complex ---
\coordinate (a) at (-2, -1.5); \coordinate (b) at (0, -0.5);
\coordinate (c) at (-1, -3); \coordinate (d) at (1.5, -2.5);
\coordinate (e) at (3, -0.8); \coordinate (f) at (3.5, -2.8);
\coordinate (g) at (-3.5, -2.8);
\fill[gdlorange!25] (a) -- (b) -- (c) -- cycle;
\fill[gdlorange!20] (b) -- (c) -- (d) -- cycle;
\fill[gdlorange!30] (b) -- (d) -- (e) -- cycle;
\draw[thick, gdlblue!60] (a)--(b) (b)--(c) (a)--(c) (c)--(d) (b)--(d)
  (b)--(e) (d)--(e) (d)--(f) (a)--(g) (c)--(g);
\foreach \p in {a,b,c,d,e,f,g} \fill[gdlblue!70] (\p) circle (3.5pt);
\node[anchor=west, font=\scriptsize] at (5, -0.9) {$\bullet$\; $0$-simplices};
\node[anchor=west, font=\scriptsize, text=gdlblue!60] at (5, -1.5)
  {\rule{0.5cm}{1.5pt}\; $1$-simplices (edges)};
\node[anchor=west, font=\scriptsize] at (5, -2.1)
  {\tikz\fill[gdlorange!25](0,0)rectangle(0.35,0.25);\; $2$-simplices (faces)};
\node[anchor=west, font=\scriptsize] at (5, -2.7) {$1$-skeleton $=G(V,E)$};
\end{tikzpicture}
\caption[Simplices and a simplicial complex]{Top: simplices of increasing
dimension --- a $0$-simplex (point), a $1$-simplex (edge), a $2$-simplex (filled
triangle) and a $3$-simplex (tetrahedron). Bottom: a simplicial complex
$\mathcal{K}$ assembled from simplices, whose $1$-skeleton forms the underlying
graph $G = (V, E)$.}
\label{fig:graphs:complex}
\end{figure}

\begin{definition}[Triangulation]
\label{def:graphs:triangulation}
A \emph{triangulation} of a compact topological manifold $M$ is a pair
$(\mathcal{K}, h)$ where $\mathcal{K}$ is a simplicial complex and
$h \colon \abs{\mathcal{K}} \to M$ is a homeomorphism from the
\emph{geometric realization} $\abs{\mathcal{K}}$ (the union of the simplices of
$\mathcal{K}$ with the induced topology) onto $M$.
\end{definition}

\begin{theorem}[Radó]
\label{thm:graphs:rado}
Every compact topological $2$-manifold admits a triangulation.
\end{theorem}

This classical result, due to \citet{rado1925aufgabe}, guarantees that every
compact surface can be decomposed into triangles compatibly with its topology ---
fundamental both for the Gauss--Bonnet theorem and for geometric deep learning on
triangulated meshes (a proof sketch appears in \Cref{ch:appendix}). Triangulability
depends crucially on dimension: in dimension $3$ it was extended by
\citet{moise1977geometric}, but in dimension $\ge 4$ there exist non-triangulable
manifolds. Simplicial complexes capture higher-order topological information ---
not only pairwise connectivity but also cycles, holes and higher-dimensional
components --- which motivates \emph{topological deep learning}
(\Cref{sec:graphs:tdl}).

\section{The topological geometry of graphs}
\label{sec:graphs:topology}

Graphs carry a rich topological and spectral structure that lets us define
discrete analogues of continuous geometric concepts. This section develops the
\emph{calculus on graphs}: combinatorial connectivity, discrete differential
operators and their spectral properties.

\subsection{Basic topological properties}
\label{sec:graphs:basic-topology}

\begin{definition}[Path in a graph]
\label{def:graphs:path}
A \emph{path} from $u$ to $v$ in a graph $G = (V, E)$ is a sequence of vertices
$u = v_0, v_1, \dots, v_k = v$ with $(v_{i-1}, v_i) \in E$ for all
$i = 1, \dots, k$. The integer $k$ is the \emph{length} of the path.
\end{definition}

\begin{definition}[Connected graph and components]
\label{def:graphs:connected}
A graph $G = (V, E)$ is \emph{connected} if for every pair of vertices
$u, v \in V$ there is a path from $u$ to $v$. A \emph{connected component} is a
maximal connected subgraph of $G$.
\end{definition}

\begin{definition}[Cycle]
\label{def:graphs:cycle}
A \emph{cycle} in $G = (V, E)$ is a closed path $v_0, v_1, \dots, v_k = v_0$ with
$k \ge 3$ and no repetition of intermediate vertices.
\end{definition}

\begin{definition}[Tree]
\label{def:graphs:tree}
A \emph{tree} is a connected graph with no cycles.
\end{definition}

These combinatorial properties encode fundamental aspects of the topology of the
discretized space. For instance, the number of connected components equals the
\emph{Betti number} $\beta_0$ of the underlying manifold, while the fundamental
cycles relate to $\beta_1$ (\Cref{sec:graphs:tdl}).

\subsection{The graph Laplacian}
\label{sec:graphs:laplacian}

The central operator of graph analysis is the \emph{Laplacian}, the discrete
counterpart of the Laplace--Beltrami operator on a manifold (\Cref{ch:geodesics}).
The labelling of the nodes carried by the matrices below is, however, an
arbitrary artefact: nothing intrinsic distinguishes ``node $1$'' from ``node
$7$''. We return to this symmetry in \Cref{sec:graphs:permutation}.

\begin{definition}[Adjacency and degree matrices]
\label{def:graphs:adjacency}
For an undirected weighted graph $G = (V, E, W)$ with $n = \abs{V}$ vertices, the
\emph{adjacency matrix} $A \in \R^{n \times n}$ has $A_{ij} = W_{ij}$ if
$(i,j) \in E$ and $A_{ij} = 0$ otherwise; for unweighted graphs
$A_{ij} \in \{0,1\}$. The \emph{degree matrix} is the diagonal matrix
$D \in \R^{n \times n}$ with $D_{ii} = \deg(i) = \sum_j A_{ij}$. The graph is
undirected when $A = A^\top$.
\end{definition}

\begin{definition}[Graph Laplacian]
\label{def:graphs:laplacian}
The (combinatorial) \emph{Laplacian} of a graph with adjacency matrix $A$ and
degree matrix $D$ is $L = D - A$, and the \emph{normalized Laplacian} (assuming
$D_{ii} > 0$) is
\[
  \Lap = D^{-1/2} L\, D^{-1/2} = I - D^{-1/2} A\, D^{-1/2}.
\]
\end{definition}

\begin{figure}[h]
\centering
\begin{tikzpicture}[font=\small,
    gnode/.style={circle, draw=gdlblue!70, fill=gdlblue!20, minimum size=20pt,
        inner sep=0pt, font=\small\bfseries}]
\node[gnode] (1) at (0,1.6) {1};
\node[gnode] (2) at (1.6,2.2) {2};
\node[gnode] (3) at (2.2,0.6) {3};
\node[gnode] (4) at (0.6,0.0) {4};
\draw[gdlgray!60, thick] (1)--(2) (1)--(4) (2)--(3) (3)--(4) (1)--(3);
\node[anchor=west, align=left] at (3.6,1.1)
  {$L = D - A =\begin{bmatrix}
     3 & -1 & -1 & -1\\
     -1 & 2 & -1 & 0\\
     -1 & -1 & 3 & -1\\
     -1 & 0 & -1 & 2
   \end{bmatrix}$};
\end{tikzpicture}
\caption[The graph Laplacian]{A small graph and its Laplacian $L = D - A$. The
diagonal records degrees, $L_{ii} = \deg(i)$, and off-diagonal entries record
adjacency, $L_{ij} = -W_{ij}$ for $(i,j) \in E$.}
\label{fig:graphs:laplacian}
\end{figure}

\begin{definition}[Positive semidefinite matrix]
\label{def:graphs:psd}
A symmetric matrix $B \in \R^{n \times n}$ is \emph{positive semidefinite}
(written $B \succeq 0$) if $x^\top B\, x \ge 0$ for all $x \in \R^n$, and
\emph{positive definite} ($B \succ 0$) if the inequality is strict for all
$x \neq 0$. By the spectral theorem, $B \succeq 0$ if and only if all its
eigenvalues are nonnegative.
\end{definition}

\begin{proposition}[Algebraic properties of the Laplacian]
\label{prop:graphs:laplacian-algebra}
The Laplacian $L = D - A$ of an undirected weighted graph satisfies:
\begin{enumerate}
  \item \textbf{Symmetry}: $L = L^\top$.
  \item \textbf{Rows sum to zero}: $L \mathbf{1} = \mathbf{0}$.
  \item \textbf{Positive semidefiniteness}: for all $x \in \R^n$,
        \[
          x^\top L\, x = \tfrac12 \sum_{(i,j) \in E} W_{ij}\,(x_i - x_j)^2 \ge 0 .
        \]
  \item \textbf{Discrete operator}: for functions $f \colon V \to \R$,
        \[
          (L f)(i) = \sum_{j \sim i} W_{ij}\,\big(f(i) - f(j)\big).
        \]
\end{enumerate}
\end{proposition}

\begin{proof}
\textbf{(1)} $D$ is diagonal and $A = A^\top$ (undirected graph), so $L = D - A$
is symmetric. \textbf{(2)} The $i$-th row of $L$ has diagonal entry
$\deg(i) = \sum_j A_{ij}$ and off-diagonal entries $-A_{ij}$, so
$\sum_j L_{ij} = \deg(i) - \sum_j A_{ij} = 0$. \textbf{(3)} Expanding,
\begin{align*}
  x^\top L\, x = x^\top D\, x - x^\top A\, x
    &= \sum_i D_{ii} x_i^2 - \sum_{i,j} A_{ij} x_i x_j \\
    &= \sum_i x_i^2 \sum_j A_{ij} - \sum_{i,j} A_{ij} x_i x_j
    = \tfrac12 \sum_{i,j} A_{ij}(x_i - x_j)^2 ,
\end{align*}
manifestly nonnegative. \textbf{(4)}
$(Lf)(i) = (Df)(i) - (Af)(i) = \sum_j A_{ij} f(i) - \sum_j A_{ij} f(j)
= \sum_{j \sim i} W_{ij}(f(i) - f(j))$.
\end{proof}

\begin{proposition}[Spectral properties of the Laplacian]
\label{prop:graphs:laplacian-props}
Let $L$ be the Laplacian of an undirected graph with $n$ vertices. Then:
\begin{enumerate}
  \item the eigenvalues are real and satisfy
        $0 = \lambda_1 \le \lambda_2 \le \cdots \le \lambda_n$, with the constant
        vector $\mathbf{1}$ an eigenvector for $\lambda_1 = 0$;
  \item the multiplicity of the eigenvalue $0$ equals the number of connected
        components of $G$.
\end{enumerate}
\end{proposition}

\begin{proof}
By \Cref{prop:graphs:laplacian-algebra}, $L$ is symmetric and positive
semidefinite, so its eigenvalues are real and nonnegative; since
$L \mathbf{1} = \mathbf{0}$, we have $\lambda_1 = 0$ with eigenvector
$\mathbf{1}$. For (2), $L x = \mathbf{0}$ iff $x^\top L\, x = 0$, which by the
quadratic form forces $(x_i - x_j)^2 = 0$ on every edge, i.e.\ $x$ is constant on
each connected component; the space of such $x$ has dimension equal to the number
of components.
\end{proof}

The quadratic form $x^\top L\, x$ measures the \emph{smoothness} of $x$ over the
graph: it is small when $x$ varies little across edges and large when $x$ changes
abruptly. This is the discrete \emph{Dirichlet energy}, and it reappears when we
analyse oversmoothing (\Cref{sec:graphs:oversmoothing}).

\paragraph{Connection with differential geometry.}
The graph Laplacian $L$ is the discrete analogue of the Laplace--Beltrami
operator $\Delta_{\mathcal{M}}$ on a Riemannian manifold (\Cref{ch:geodesics}),
which acts on smooth functions $f \colon \mathcal{M} \to \R$ by
$\Delta_{\mathcal{M}} f = \dive(\grad f)$. The analogy is precise: on a continuous
manifold $\Delta_{\mathcal{M}} f$ measures the divergence of the gradient, while
on a discrete graph $(L f)(i) = \sum_{j \sim i} (f(i) - f(j))$ measures the
difference between $f(i)$ and the average of its neighbours. Both generalize the
Euclidean Laplacian $\nabla^2 f$ to non-Euclidean domains. We make the limiting
correspondence rigorous in \Cref{sec:graphs:convergence}.

\section{Topological deep learning}
\label{sec:graphs:tdl}

The Laplacian $L = D - A$ acts exclusively on node signals ($0$-chains). Yet the
simplicial complexes of \Cref{def:graphs:complex} carry a far richer algebraic
structure: signals defined on edges, triangles and higher-order simplices. To
quantify it we import the tools of algebraic topology.

\begin{definition}[Chains]
\label{def:graphs:chains}
Let $\mathcal{K}$ be a simplicial complex. A \emph{$k$-chain} is a formal linear
combination of $k$-simplices with coefficients in a field (typically $\R$ or
$\Z_2$):
\[
  C_k(\mathcal{K}) = \Big\{ \textstyle\sum_i \alpha_i \sigma_i^{(k)}
    : \alpha_i \in \R,\ \sigma_i^{(k)} \in \mathcal{K} \text{ a $k$-simplex} \Big\}.
\]
\end{definition}

\begin{definition}[Boundary operator]
\label{def:graphs:boundary}
The \emph{boundary operator}
$\partial_k \colon C_k(\mathcal{K}) \to C_{k-1}(\mathcal{K})$ is defined on an
oriented $k$-simplex $[v_0, v_1, \dots, v_k]$ by
\[
  \partial_k [v_0, v_1, \dots, v_k]
    = \sum_{i=0}^{k} (-1)^i\, [v_0, \dots, \widehat{v_i}, \dots, v_k],
\]
where $\widehat{v_i}$ denotes omission of the vertex $v_i$.
\end{definition}

\begin{proposition}[Fundamental property of the boundary]
\label{prop:graphs:boundary-squared}
The boundary operator satisfies $\partial_{k-1} \circ \partial_k = 0$ for all
$k$: the boundary of a boundary is empty.
\end{proposition}

\begin{proof}
It suffices to verify the identity on an oriented $k$-simplex
$[v_0, \dots, v_k]$:
\[
  \partial_{k-1}\big(\partial_k [v_0, \dots, v_k]\big)
    = \sum_{i=0}^{k} (-1)^i \sum_{j \neq i} (-1)^{j'}
      [v_0, \dots, \widehat{v_i}, \dots, \widehat{v_j}, \dots, v_k],
\]
where $j' = j$ if $j < i$ and $j' = j-1$ if $j > i$. Each codimension-$2$ face
obtained by omitting two vertices $v_i, v_j$ with $i < j$ appears exactly twice:
once with sign $(-1)^{i+j}$ (omitting $v_i$ first, then $v_j$) and once with sign
$(-1)^{i+j-1}$ (omitting $v_j$ first). Since
$(-1)^{i+j} + (-1)^{i+j-1} = 0$, all terms cancel.
\end{proof}

\begin{definition}[Homology groups]
\label{def:graphs:homology}
The \emph{homology groups} of a simplicial complex $\mathcal{K}$ are
\[
  H_k(\mathcal{K}) = \ker(\partial_k) / \im(\partial_{k+1})
    = Z_k(\mathcal{K}) / B_k(\mathcal{K}),
\]
where $Z_k = \ker(\partial_k)$ are the \emph{$k$-cycles} (chains with no
boundary) and $B_k = \im(\partial_{k+1})$ are the \emph{$k$-boundaries} (chains
that are the boundary of something).
\end{definition}

\begin{definition}[Betti numbers]
\label{def:graphs:betti}
The \emph{$k$-th Betti number} of $\mathcal{K}$ is $\beta_k = \dim H_k(\mathcal{K})$,
which counts: $\beta_0$, the number of connected components; $\beta_1$, the
number of independent ``holes'' or cycles; $\beta_2$, the number of
three-dimensional cavities.
\end{definition}

\begin{example}[Betti numbers of surfaces]
\label{ex:graphs:betti}
For the sphere $S^2$: $\beta_0 = 1$ (connected), $\beta_1 = 0$ (no holes),
$\beta_2 = 1$ (one cavity). For the torus $T^2$: $\beta_0 = 1$, $\beta_1 = 2$
(two independent cycles), $\beta_2 = 1$. For the Klein bottle: $\beta_0 = 1$,
$\beta_1 = 1$ (over $\Z_2$).
\end{example}

Topological deep learning extends message passing (\Cref{def:graphs:mpnn}) and
graph analysis to these higher-order structures through the \emph{Hodge
Laplacian}, which generalizes $L$ to arbitrary $k$-chains.

\subsection{The Hodge Laplacian}
\label{sec:graphs:hodge}

The boundary operator and the identity $\partial_{k-1} \circ \partial_k = 0$ give
the \emph{chain complex}
\[
  \cdots \xrightarrow{\partial_{k+2}} C_{k+1}(\mathcal{K})
    \xrightarrow{\partial_{k+1}} C_k(\mathcal{K})
    \xrightarrow{\partial_k} C_{k-1}(\mathcal{K})
    \xrightarrow{\partial_{k-1}} \cdots .
\]

\begin{definition}[Coboundary operator]
\label{def:graphs:coboundary}
Equip each chain space $C_k(\mathcal{K})$ with the standard inner product in
which the oriented $k$-simplices form an orthonormal basis. The
\emph{coboundary operator}
$\partial_k^\top \colon C_{k-1}(\mathcal{K}) \to C_k(\mathcal{K})$ is the formal
adjoint of $\partial_k$, the unique linear operator with
$\inner{\partial_k^\top c}{\sigma} = \inner{c}{\partial_k \sigma}$ for all
$c \in C_{k-1}$, $\sigma \in C_k$. Combinatorially, for a $(k{-}1)$-simplex
$\tau$ it aggregates all $k$-simplices containing $\tau$ as a face, with the
inherited incidence signs. Transposing
$\partial_k \circ \partial_{k+1} = 0$ gives
$\partial_{k+1}^\top \circ \partial_k^\top = 0$.
\end{definition}

\begin{definition}[Hodge Laplacian]
\label{def:graphs:hodge}
For a finite simplicial complex $\mathcal{K}$ with boundary operators
$\partial_k$, the \emph{Hodge Laplacian of order $k$} is the operator
$L_k \colon C_k(\mathcal{K}) \to C_k(\mathcal{K})$,
\[
  L_k = \underbrace{\partial_k^\top \partial_k}_{L_k^{\mathrm{down}}}
      + \underbrace{\partial_{k+1} \partial_{k+1}^\top}_{L_k^{\mathrm{up}}},
\]
where $L_k^{\mathrm{down}}$ captures \emph{lower} adjacency (two $k$-simplices are
adjacent if they share a $(k{-}1)$-face) and $L_k^{\mathrm{up}}$ captures
\emph{upper} adjacency (two $k$-simplices are adjacent if they belong to a common
$(k{+}1)$-simplex).
\end{definition}

\begin{proposition}[Properties of the Hodge Laplacian]
\label{prop:graphs:hodge-props}
The Hodge Laplacian $L_k$ satisfies:
\begin{enumerate}
  \item \textbf{Symmetry and positive semidefiniteness}: $L_k$ is symmetric and
        $\inner{x}{L_k x} = \norm{\partial_k x}^2 + \norm{\partial_{k+1}^\top x}^2
        \ge 0$ for every $k$-chain $x$.
  \item \textbf{Recovery of the graph Laplacian}: for $k = 0$, with
        $\partial_0 = 0$ by convention, $L_0 = \partial_1 \partial_1^\top$
        coincides with $L = D - A$ of \Cref{def:graphs:laplacian}.
  \item \textbf{Connection with homology}: $\ker(L_k) \cong H_k(\mathcal{K})$,
        i.e.\ the harmonic $k$-chains correspond to the $k$-th homology group.
\end{enumerate}
\end{proposition}

\begin{proof}
\textbf{(1)} $L_k$ is a sum of Gram matrices, hence symmetric; and
$\inner{x}{L_k x} = \inner{x}{\partial_k^\top \partial_k x}
+ \inner{x}{\partial_{k+1}\partial_{k+1}^\top x}
= \norm{\partial_k x}^2 + \norm{\partial_{k+1}^\top x}^2 \ge 0$.
\textbf{(2)} For $k = 0$, $C_{-1}(\mathcal{K})$ is trivial so $\partial_0 = 0$
and $L_0 = \partial_1 \partial_1^\top$. If $\partial_1$ is the oriented incidence
matrix, then $\partial_1 \partial_1^\top$ has $(i,j)$ entry $\deg(i)$ when
$i = j$, $-1$ when $(i,j) \in E$, and $0$ otherwise; that is,
$\partial_1 \partial_1^\top = D - A = L$. \textbf{(3)} By (1), $L_k x = 0$ iff
$\partial_k x = 0$ and $\partial_{k+1}^\top x = 0$ simultaneously: the first says
$x \in Z_k$, the second that $x \perp \im(\partial_{k+1})$. Thus $x$ is a cycle
orthogonal to all boundaries, and the class $[x] \in Z_k / B_k = H_k(\mathcal{K})$
gives the isomorphism $\ker(L_k) \cong H_k(\mathcal{K})$.
\end{proof}

The central result of discrete Hodge theory provides an orthogonal decomposition
of the chain space that generalizes graph Fourier analysis.

\begin{theorem}[Hodge decomposition]
\label{thm:graphs:hodge-decomp}
For a finite simplicial complex $\mathcal{K}$, the space of $k$-chains admits the
orthogonal decomposition
\[
  C_k(\mathcal{K})
    = \underbrace{\im(\partial_{k+1})}_{\text{boundaries (``gradients'')}}
    \oplus \underbrace{\ker(L_k)}_{\text{harmonics}}
    \oplus \underbrace{\im(\partial_k^\top)}_{\text{coboundaries (``curls'')}} .
\]
The three components capture complementary aspects of a simplicial signal:
boundaries represent irrotational flow, harmonics encode the global topology
($H_k$), and coboundaries represent pure circulation.
\end{theorem}

\begin{proof}
\textbf{Mutual orthogonality.} For $x = \partial_{k+1} a \in \im(\partial_{k+1})$
and $y = \partial_k^\top b \in \im(\partial_k^\top)$, using
$\partial_k \partial_{k+1} = 0$,
$\inner{x}{y} = \inner{\partial_{k+1} a}{\partial_k^\top b}
= \inner{\partial_k \partial_{k+1} a}{b} = 0$. For $z \in \ker(L_k)$ we have
$\partial_k z = 0$ and $\partial_{k+1}^\top z = 0$, hence
$\inner{z}{\partial_{k+1} a} = \inner{\partial_{k+1}^\top z}{a} = 0$ and
$\inner{z}{\partial_k^\top b} = \inner{\partial_k z}{b} = 0$.
\textbf{Completeness.} In finite dimension it suffices to show
$(V_1 \oplus V_2)^\perp = V_3$ with $V_1 = \im(\partial_{k+1})$,
$V_2 = \im(\partial_k^\top)$, $V_3 = \ker(L_k)$. If
$w \perp V_1$ and $w \perp V_2$, then $\partial_{k+1}^\top w = 0$ and
$\partial_k w = 0$, so
$L_k w = \partial_k^\top \partial_k w + \partial_{k+1}\partial_{k+1}^\top w = 0$,
i.e.\ $w \in V_3$.
\end{proof}

\begin{example}[Hodge Laplacian $L_1$ of a triangular complex]
\label{ex:graphs:hodge-L1}
Consider a complex $\mathcal{K}$ with vertices $\{a, b, c, d\}$, edges
$e_1 = [a,b]$, $e_2 = [b,c]$, $e_3 = [a,c]$, $e_4 = [c,d]$, and one $2$-simplex
$\sigma = [a,b,c]$. From $\partial_2[a,b,c] = e_2 - e_3 + e_1$,
\[
  \partial_2 = \begin{pmatrix} 1 \\ 1 \\ -1 \\ 0 \end{pmatrix}, \qquad
  \partial_1 = \begin{pmatrix}
    -1 & 0 & -1 & 0 \\ 1 & -1 & 0 & 0 \\ 0 & 1 & 1 & -1 \\ 0 & 0 & 0 & 1
  \end{pmatrix},
\]
and one verifies $\partial_1 \partial_2 = 0$. Then
\[
  L_1^{\mathrm{down}} = \partial_1^\top \partial_1 = \begin{pmatrix}
    2 & -1 & 1 & 0 \\ -1 & 2 & 1 & -1 \\ 1 & 1 & 2 & -1 \\ 0 & -1 & -1 & 2
  \end{pmatrix}, \quad
  L_1^{\mathrm{up}} = \partial_2 \partial_2^\top = \begin{pmatrix}
    1 & 1 & -1 & 0 \\ 1 & 1 & -1 & 0 \\ -1 & -1 & 1 & 0 \\ 0 & 0 & 0 & 0
  \end{pmatrix}.
\]
Here $L_1^{\mathrm{down}}$ connects edges sharing a vertex (lower adjacency),
while $L_1^{\mathrm{up}}$ connects edges coexisting in $\sigma$ (upper
adjacency). The edge $e_4$, belonging to no triangle, has a zero row and column
in $L_1^{\mathrm{up}}$.
\end{example}

\subsection{Simplicial message passing}
\label{sec:graphs:simplicial-mp}

The decomposition $L_k = L_k^{\mathrm{down}} + L_k^{\mathrm{up}}$ reveals that
each $k$-simplex receives information from two geometrically distinct sources:
its $(k{-}1)$-faces (lower adjacency) and the $(k{+}1)$-simplices that contain it
(upper adjacency). This dual structure motivates a natural generalization of
graph message passing to simplicial complexes: a target $k$-simplex $\sigma$
updates its state from lower messages $m^{\downarrow}$ (from its boundary), upper
messages $m^{\uparrow}$ (from its coboundary), and horizontal messages
$m^{\leftrightarrow}$ (from adjacent simplices of the same dimension), as
illustrated in \Cref{fig:graphs:simplicial-mp}.

\begin{figure}[h]
\centering
\begin{tikzpicture}[>=Stealth, font=\small,
    msg/.style={line width=1.6pt, -{Stealth[length=6pt, width=4pt]}}]
\coordinate (A) at (0, 0); \coordinate (B) at (4.3, 0);
\coordinate (C) at (2.15, 3.0);
\coordinate (Dd) at (-3.0, 1.0); \coordinate (Ee) at (-2.6, -1.7);
\coordinate (Ff) at (7.3, 1.0);
\fill[gdlorange!15] (A) -- (B) -- (C) -- cycle;
\draw[gdlorange!40, line width=0.8pt] (A) -- (B) -- (C) -- cycle;
\node[gdlorange!70!black] at (2.15, 1.35) {$\tau^2$};
\draw[gdlgray!50, line width=1pt] (A) -- (C); \draw[gdlgray!50, line width=1pt] (B) -- (C);
\draw[gdlpurple!30, line width=2pt] (A) -- (Dd);
\draw[gdlpurple!30, line width=2pt] (A) -- (Ee);
\draw[gdlpurple!30, line width=2pt] (B) -- (Ff);
\draw[gdlred!70, line width=3.5pt] (A) -- (B);
\node[gdlred!80!black, font=\bfseries] at (2.15, -0.5) {$\sigma^1$ (target)};
\foreach \v in {A,B,C,Dd,Ee,Ff} {\fill[gdlblue!60] (\v) circle (4pt);
  \draw[gdlblue!80] (\v) circle (4pt);}
% lower (green) from vertices
\draw[msg, gdlgreen!70] ($(A)+(-0.3,-1.6)$) to[bend right=25] ($(A)+(0.45,-0.22)$);
\draw[msg, gdlgreen!70] ($(B)+(0.3,-1.6)$) to[bend left=25] ($(B)+(-0.45,-0.22)$);
\node[gdlgreen!70!black, font=\scriptsize, fill=white, fill opacity=0.85,
  text opacity=1, inner sep=2pt] at (2.15, -1.6) {$m^{\downarrow}$ (boundary)};
% upper (orange) from triangle
\draw[msg, gdlorange!70] (2.15, 2.4) -- (2.15, 0.28);
\node[gdlorange!70!black, font=\scriptsize, anchor=south west] at (2.6, 1.95)
  {$m^{\uparrow}$ (coboundary)};
% horizontal (purple)
\draw[msg, gdlpurple!70] ($(A)!0.5!(Dd)$) to[bend right=20] ($(A)+(0.55,0.3)$);
\draw[msg, gdlpurple!70] ($(A)!0.5!(Ee)$) to[bend left=20] ($(A)+(0.55,-0.3)$);
\draw[msg, gdlpurple!70] ($(B)!0.5!(Ff)$) to[bend left=20] ($(B)+(-0.55,0.3)$);
\node[gdlpurple!70!black, font=\scriptsize, anchor=east, fill=white,
  fill opacity=0.85, text opacity=1, inner sep=2pt] at (-3.1, 0)
  {$m^{\leftrightarrow}$ (adjacency)};
\end{tikzpicture}
\caption[Simplicial message passing]{Message passing on a simplicial complex for
a target edge $\sigma^1$ (red). Lower messages $m^{\downarrow}$ arrive from its
boundary vertices; upper messages $m^{\uparrow}$ from the coboundary triangle
$\tau^2$ containing it; horizontal messages $m^{\leftrightarrow}$ from adjacent
edges sharing a vertex.}
\label{fig:graphs:simplicial-mp}
\end{figure}

\paragraph{Architectural variants and extensions.}
The topological deep learning framework admits several specializations:
\emph{topological convolutions} (spectral filters on the Hodge Laplacian,
parametrized by polynomials of $L_k$); \emph{simplicial attention}, which weights
simplicial messages by orientation and topological relevance; \emph{cell
complexes}, generalizing simplicial complexes to cells of arbitrary shape;
\emph{simplicial diffusion}, extending the heat equation to $L_k$; and
higher-order attention networks combining multi-head attention with simplicial
topology. These are developed further in \Cref{ch:gauges,ch:applications}.

\section{Convergence of graphs to manifolds}
\label{sec:graphs:convergence}

The question that ties graph theory to differential geometry is: under what
conditions does a discrete graph faithfully approximate a continuous manifold?
The following theorem --- one of the central results of manifold learning ---
combines tools from differential geometry, random graph theory and asymptotic
analysis to give a rigorous answer.

\begin{theorem}[Convergence of graphs to manifolds]
\label{thm:graphs:convergence}
Let $(\mathcal{M}, g)$ be a compact Riemannian manifold of dimension $d$ embedded
in $\R^N$, and let $\{x_1, \dots, x_n\}$ be points sampled on $\mathcal{M}$
according to a density $p(x) \ge p_{\min} > 0$. Build a graph $G_n$ with weights
$W_{ij} = k_\epsilon(\norm{x_i - x_j})$ where
$k_\epsilon(t) = \epsilon^{-d}\, k(t/\epsilon)$ and $k \colon \R_+ \to \R_+$ is a
smooth, compactly supported radial kernel. Then, under the scaling conditions
$\epsilon = \epsilon_n \to 0$ and $n\epsilon^{d+2} \to \infty$, the
density-normalized graph Laplacian converges to the weighted Laplace--Beltrami
operator:
\[
  \lim_{n \to \infty} \frac{1}{p(x_i)\, \epsilon^2}\, L_n f(x_i)
    = \Delta_{\mathcal{M}} f(x_i)
      + \inner{\nabla \log p(x_i)}{\nabla f(x_i)},
\]
in probability, uniformly for $f \in C^2(\mathcal{M})$.
\end{theorem}

\begin{proof}[Proof sketch]
We outline the seven steps combining differential geometry, functional analysis
and probability.

\textbf{Step 1 (Normalization).} Define
$L_n f(x_i) = \frac{1}{n\epsilon^d} \sum_{j} k_\epsilon(\norm{x_i - x_j})
[f(x_j) - f(x_i)]$; the factor $n\epsilon^d$ represents the expected local
neighbourhood volume.

\textbf{Step 2 (Normal coordinates).} Work in geodesic normal coordinates
centred at $x_i$, where $g_{jk}(y) = \delta_{jk} + O(\norm{y}^2)$, so the
manifold is locally approximated by its tangent space
$T_{x_i}\mathcal{M} \cong \R^d$. For nearby $x_j$ write
$x_j = \exp_{x_i}(\epsilon v_j)$ with $v_j \in T_{x_i}\mathcal{M}$.

\textbf{Step 3 (Taylor expansion).}
$f(x_j) - f(x_i) = \epsilon \inner{\nabla f(x_i)}{v_j}
+ \tfrac{\epsilon^2}{2} \Hess_f(x_i)(v_j, v_j) + O(\epsilon^3)$.

\textbf{Steps 4--5 (Asymptotic average).} Substituting and applying the strong
law of large numbers, $\frac1n \sum_j F(x_j) \to \int_{\mathcal{M}} F\, p\, dV_g$,
and changing variables $y = \exp_{x_i}(\epsilon v)$ while expanding the density
$p(\exp_{x_i}(\epsilon v)) = p(x_i) + \epsilon \inner{\nabla p(x_i)}{v} + O(\epsilon^2)$.

\textbf{Step 6 (First-order term).} The term
$p(x_i) \int k(\norm{v}) \inner{\nabla f}{v}\, dv = 0$ vanishes by radial
symmetry, but the cross term of order $\epsilon$ survives: using
$\int k(\norm{v}) v_\alpha v_\beta\, dv = \frac1d \big(\int k(\norm{v})\norm{v}^2 dv\big)
\delta_{\alpha\beta}$, it contributes proportionally to
$\inner{\nabla f(x_i)}{\nabla p(x_i)}$.

\textbf{Step 7 (Second-order term).} Using
$\tr(\Hess_f) = \Delta_{\mathcal{M}} f$ and radial symmetry, the second-order term
gives $\frac{p(x_i)}{2d}\big(\int k(\norm{v})\norm{v}^2 dv\big)
\Delta_{\mathcal{M}} f(x_i)$. Choosing $k$ so that
$\int_{\R^d} k(\norm{v})\norm{v}^2 dv = 2d$ and combining Steps 6--7,
\[
  \lim_{n\to\infty} \frac{1}{\epsilon^2} L_n f(x_i)
    = p(x_i)\Delta_{\mathcal{M}} f(x_i) + \inner{\nabla p(x_i)}{\nabla f(x_i)} ,
\]
and dividing by $p(x_i)$ with $\nabla p / p = \nabla \log p$ yields the claim.
\end{proof}

The term $\inner{\nabla \log p}{\nabla f}$ reflects the influence of nonuniform
sampling and can be removed by more sophisticated Laplacian constructions
\citep{coifman2006diffusion, hein2007graph}. Convergence requires
$\epsilon \to 0$ (vanishing local radius) yet $n\epsilon^{d+2} \to \infty$
(enough points per neighbourhood); the optimal rate is
$\epsilon \sim (\log n / n)^{1/(d+2+\alpha)}$ for small $\alpha > 0$, balancing
approximation bias against sampling variance.

\begin{remark}[Convergence rates and scaling regimes]
\label{rem:graphs:convergence-rates}
The quality of convergence depends critically on the relation between $n$ and
$\epsilon$. Three regimes arise: \emph{underconnected} ($n\epsilon^d \to 0$), the
graph is disconnected and there is no convergence; \emph{optimal}
($n\epsilon^d \sim \log n$), convergence at rate
$O(\epsilon^2 + 1/\sqrt{n\epsilon^d})$; and \emph{overconnected}
($n\epsilon^d \to \infty$ too fast), convergence but with loss of locality. For
$n = 10\,000$ points uniform on $S^2$ ($d=2$), the error in the first nontrivial
eigenvalue is $\approx 15\%$ for $\epsilon = 0.05$ (underconnected),
$\approx 2\%$ for $\epsilon = 0.15$ (optimal), and $\approx 5\%$ for
$\epsilon = 0.40$ (overconnected). The optimal
$\epsilon^* \approx (\log n / n)^{1/4} \approx 0.14$ matches the theory.
\end{remark}

This result has profound consequences for geometric deep learning. It gives a
\emph{theoretical justification} for graph neural networks on manifold-valued
data; it \emph{guides graph construction} (how to choose $\epsilon$ as a function
of $n$ and the intrinsic dimension $d$); it exposes \emph{fundamental
limitations} (finite data yields only approximations, and the smooth geometry
between sampled points is lost); and it explains a \emph{computational trade-off}
--- graphs are the practical middle ground between Euclidean representations
(which ignore geometry) and continuous manifolds (elegant but intractable).
Spectral graph networks such as ChebNet and the graph convolutional network
inherit geometric properties of the continuous domain precisely because they
approximate spectral filters of the Laplacian.

\section{Spectral analysis on graphs}
\label{sec:graphs:spectral}

Spectral graph analysis extends classical Fourier analysis to the discrete
domain (the harmonic-analysis foundations are developed in \Cref{ch:grids}). The
eigenvectors of the Laplacian form an orthogonal basis analogous to the sinusoids
of Euclidean Fourier analysis, letting us decompose functions into
``frequencies'' over the graph.

\subsection{The graph Fourier transform}
\label{sec:graphs:gft}

Because $\Lap$ is symmetric it admits an orthonormal eigendecomposition
$\Lap = U \Lambda U^\top$ with $\Lambda = \operatorname{diag}(\lambda_1, \dots,
\lambda_n)$, $0 = \lambda_1 \le \lambda_2 \le \cdots \le \lambda_n \le 2$, and
$U = [u_1, \dots, u_n]$ orthonormal.

\begin{definition}[Graph Fourier transform]
\label{def:graphs:gft}
The \emph{graph Fourier transform} of a signal $f \colon V \to \R$ (viewed as
$f \in \R^n$) is $\hat f = U^\top f \in \R^n$, with inverse
$f = U \hat f = \sum_i \hat f_i\, u_i$. The coefficient
$\hat f_i = \inner{f}{u_i}$ is the amplitude of $f$ at frequency $\lambda_i$.
\end{definition}

The eigenvectors $u_i$ play the role of the complex exponentials
$e^{2\pi \imath \omega \cdot x}$, and the eigenvalues $\lambda_i$ the role of the
frequencies $\omega$. \Cref{fig:graphs:eigenvectors} illustrates how small
eigenvalues correspond to smooth (low-frequency) eigenvectors and large
eigenvalues to oscillating (high-frequency) ones.

\begin{figure}[h]
\centering
\begin{tikzpicture}[font=\small]
\foreach \k/\lab/\xs in {0/{$\lambda$ small}/0, 1/{$\lambda$ medium}/4.6,
                         2/{$\lambda$ large}/9.2}{
  \begin{scope}[shift={(\xs,0)}]
    \begin{scope}[on background layer]
      \fill[gdlgray!8, rounded corners=2pt] (-0.4,-0.4) rectangle (3.4,2.7);
    \end{scope}
    \foreach \i in {0,1,2,3} \foreach \j in {0,1,2}
      \coordinate (n\k-\i-\j) at (\i,\j);
    \foreach \i in {0,1,2,3} \foreach \j in {0,1,2}{
      \pgfmathsetmacro{\val}{ifthenelse(\k==0, 0.6,
        ifthenelse(\k==1, 0.6*cos(180*\i), 0.6*cos(180*\i)*cos(180*\j)))}
      \pgfmathsetmacro{\cc}{50*\val+50}
      \ifdim\val pt>0pt
        \fill[gdlblue!\cc] (\i,\j) circle (5pt);
      \else
        \pgfmathsetmacro{\dd}{-100*\val}
        \fill[gdlred!\dd] (\i,\j) circle (5pt);
      \fi
      \draw[gdlgray!50] (\i,\j) circle (5pt);
    }
    \node[font=\scriptsize\bfseries] at (1.5,-0.85) {\lab};
  \end{scope}
}
\end{tikzpicture}
\caption[Laplacian eigenvectors]{Laplacian eigenvectors on a grid graph. Small
eigenvalues correspond to smooth functions (low frequency, nodes vary little
between neighbours); large eigenvalues correspond to oscillatory functions (high
frequency, sign changes between adjacent nodes). Blue and red denote the sign of
the eigenvector entry.}
\label{fig:graphs:eigenvectors}
\end{figure}

\begin{theorem}[Frequency interpretation of eigenvectors]
\label{thm:graphs:frequency}
The Laplacian eigenvectors $\{u_i\}$ satisfy:
\begin{enumerate}
  \item \textbf{Small eigenvalues $\Leftrightarrow$ smooth functions}: if
        $\Lap u_i = \lambda_i u_i$ with $\lambda_i$ small, then
        $u_i^\top \Lap\, u_i = \lambda_i$ is small, which forces $u_i$ to vary
        smoothly across edges.
  \item \textbf{Large eigenvalues $\Leftrightarrow$ oscillatory functions}: if
        $\lambda_i \approx 2$, then $u_i$ oscillates rapidly, changing sign
        between adjacent nodes.
  \item \textbf{Analogy with frequencies}: $\lambda_i$ plays the role of
        ``frequency squared'', $\lambda_i \sim \omega_i^2$.
\end{enumerate}
\end{theorem}

\begin{proof}
\textbf{(1)} For the normalized Laplacian,
$u_i^\top \Lap\, u_i = \tfrac12 \sum_{(j,k)\in E} W_{jk}
\big(u_i(j)/\sqrt{d_j} - u_i(k)/\sqrt{d_k}\big)^2$; if $\lambda_i$ is small this
sum is small, requiring $u_i$ to vary little across edges.
\textbf{(2)} Since $\lambda_{\max} \le 2$, values near $2$ are attained when
$u_i$ changes sign frequently between neighbours, maximizing
$(u_i(j) - u_i(k))^2$. \textbf{(3)} The Euclidean Laplacian has eigenfunctions
$e^{\imath \omega \cdot x}$ with eigenvalues $-\omega^2$; with the sign
convention, $\lambda \sim \omega^2$.
\end{proof}

\subsection{Spectral filters}
\label{sec:graphs:filters}

By analogy with the classical convolution theorem --- convolution in space is
multiplication in frequency --- one defines convolution on a graph as
multiplication of Fourier coefficients by a filter.

\begin{definition}[Spectral graph filter]
\label{def:graphs:spectral-filter}
A \emph{spectral filter} with transfer function $g_\theta \colon \R \to \R$ acts
on a signal $f$ by
\[
  g_\theta \star f
    = U\, g_\theta(\Lambda)\, U^\top f
    = U\, \operatorname{diag}\!\big(g_\theta(\lambda_1), \dots, g_\theta(\lambda_n)\big)\, U^\top f .
\]
\end{definition}

This is the foundation of \emph{spectral} graph networks~\citep%
{bruna2014spectralnetworkslocallyconnected}, which learn the transfer function
$g_\theta(\lambda)$ directly. The analogy with continuous harmonic analysis is
exact in the limit: $\Delta_{\mathcal{M}}$ corresponds to $L$, its eigenfunctions
$\Delta_{\mathcal{M}}\phi_i = \lambda_i \phi_i$ to the eigenvectors
$L u_i = \lambda_i u_i$, the manifold Fourier transform to $\hat f = U^\top f$,
and smooth (resp.\ oscillatory) functions to small (resp.\ large) eigenvalues.

Attractive as it is, the naive spectral approach has three drawbacks. First, it
requires the full eigendecomposition of $\Lap$, an $O(n^3)$ computation. Second,
a freely parametrized $g_\theta(\lambda)$ uses $O(n)$ parameters and is not
localized in space. Third --- and most fundamentally --- the eigenbasis $U$ is
tied to one graph, so a filter learned on one graph does not transfer to another,
breaking the inductive setting in which graphs vary from example to example. The
polynomial filters of \Cref{sec:graphs:cheb} remove all three drawbacks at once;
before that, we develop the spatial viewpoint that they meet.

\section{Permutation symmetry}
\label{sec:graphs:permutation}

The matrix representation of a graph carries an arbitrary artefact: the labelling
of the nodes. Renumbering the atoms of a molecule does not change the molecule. A
permutation of the node labels is therefore a symmetry of the data, and any model
we build must respect it. A permutation acts simultaneously on the rows of $X$
and on both axes of $A$.

\begin{definition}[Permutation action]
\label{def:graphs:permutation-action}
Let $\group_n$ be the symmetric group of permutations of $\{1, \dots, n\}$,
represented by permutation matrices $P \in \R^{n \times n}$ (one entry $1$ per
row and column, zero elsewhere). A permutation $P$ acts on node features and on
the adjacency matrix by
\[
  X \mapsto P X, \qquad A \mapsto P A P^\top .
\]
\end{definition}

We want a node-level computation $F(X, A)$ to commute with relabelling, and a
graph-level computation $f(X, A)$ to ignore relabelling entirely. These are the
two faces of symmetry from \Cref{ch:priors}, specialized to $\group_n$.

\begin{definition}[Permutation equivariance and invariance]
\label{def:graphs:equiv-inv}
A node-wise function $F \colon \R^{n \times d} \times \R^{n \times n} \to
\R^{n \times d'}$ is \emph{permutation equivariant} if
$F(P X, P A P^\top) = P\, F(X, A)$ for all $P \in \group_n$. A graph-wise
function $f \colon \R^{n \times d} \times \R^{n \times n} \to \R^{d'}$ is
\emph{permutation invariant} if $f(P X, P A P^\top) = f(X, A)$ for all
$P \in \group_n$.
\end{definition}

These conditions parallel the equivariance and invariance demanded of
convolutional and group-equivariant layers, but the relevant group is now the
discrete, finite symmetric group rather than translations or rotations. The
blueprint of \Cref{ch:priors} applies verbatim: stack equivariant layers, which
compose into an equivariant map, and append an invariant readout (a
permutation-invariant pooling such as sum or mean over nodes) only at the end,
when a single graph-level prediction is required.

The constraint is severe. The only linear maps on node features that are
permutation equivariant for \emph{every} graph are spanned by two operators: the
identity (each node transforms itself) and the all-ones aggregation (each node
sums the features of all nodes). To obtain a richer, structure-aware family we
must let the aggregation respect the edges: each node may combine its own feature
with an aggregate of its neighbours' features, where the aggregate is itself
invariant to the order in which neighbours are listed. This single observation
--- aggregate over a neighbourhood with a permutation-invariant operator, then
update --- is the seed of every graph neural network, and we turn to it next.

\section{Message passing}
\label{sec:graphs:message-passing}

A graph neural network is built by repeating a local, permutation-equivariant
update in which every node refreshes its state from the states of its neighbours.
\citet{gilmer2017neural} crystallized this idea into the \emph{message passing}
framework, which has since become the lingua franca of learning on graphs.

\begin{definition}[Message passing neural network]
\label{def:graphs:mpnn}
A \emph{message passing neural network} (MPNN) operates on a graph $G = (V, E)$
with node features $\{h_v^{(\ell)}\}$ and optional edge features $\{e_{vu}\}$.
Layer $\ell$ updates each node in two phases:
\begin{align}
  m_v^{(\ell+1)} &= \bigoplus_{u \in \Nbr(v)}
      M_\ell\!\left(h_v^{(\ell)}, h_u^{(\ell)}, e_{vu}\right)
      && \text{(message and aggregation)}, \\
  h_v^{(\ell+1)} &= U_\ell\!\left(h_v^{(\ell)}, m_v^{(\ell+1)}\right)
      && \text{(update)},
\end{align}
where $M_\ell$ is a learnable message function, $U_\ell$ a learnable update
function, and $\bigoplus$ a permutation-invariant aggregator (sum, mean, or
maximum).
\end{definition}

The aggregator is the load-bearing component: because $\bigoplus$ does not depend
on the order of its arguments, the layer is permutation equivariant in the sense
of \Cref{def:graphs:equiv-inv}, whatever the message and update functions. After
$\ell$ layers each node has gathered information from its $\ell$-hop
neighbourhood, exactly as stacking convolutional layers grows the receptive
field. A graph-level prediction is produced by a final invariant pooling,
$h_G = \bigoplus_{v \in V} h_v^{(L)}$, followed by a standard predictor.
\Cref{fig:graphs:message-passing} depicts the three phases for a single node.

\begin{figure}[h]
\centering
\begin{tikzpicture}[>=Stealth, font=\small,
    centernode/.style={circle, draw=gdlred!70, fill=gdlred!20,
        minimum size=22pt, inner sep=0pt, font=\small\bfseries, line width=1.2pt},
    neighnode/.style={circle, draw=gdlblue!60, fill=gdlblue!20,
        minimum size=20pt, inner sep=0pt, font=\small}]
\node[font=\bfseries, text=gdlorange!70!black] at (1.5, 4.3) {1. Message};
\node[font=\bfseries, text=gdlpurple!70!black] at (6.6, 4.3) {2. Aggregate};
\node[font=\bfseries, text=gdlgreen!70!black] at (10.6, 4.3) {3. Update};
\begin{scope}
  \node[centernode] (vi) at (1.5, 2.0) {$v$};
  \node[neighnode] (v1) at (0.0, 3.5) {$u_1$};
  \node[neighnode] (v2) at (3.0, 3.5) {$u_2$};
  \node[neighnode] (v3) at (3.1, 1.0) {$u_3$};
  \node[neighnode] (v4) at (0.0, 0.5) {$u_4$};
  \foreach \n in {v1,v2,v3,v4}
    \draw[->, line width=1.4pt, gdlorange!70, shorten >=5pt, shorten <=5pt] (\n) -- (vi);
  \node[font=\footnotesize, text=gdlorange!70!black, anchor=north, align=center] at (1.5, -0.2)
    {$m_{u\to v}=M_\ell(h_v,h_u,e_{vu})$};
\end{scope}
\draw[->, line width=1.6pt, gdlgray!70] (4.2, 2.0) -- (5.0, 2.0);
\begin{scope}[shift={(5.2,0)}]
  \node[circle, draw=gdlpurple!70, fill=gdlpurple!15, minimum size=38pt,
        inner sep=0pt, line width=1.2pt, font=\large] (agg) at (1.4, 2.0) {$\bigoplus$};
  \node[font=\footnotesize, text=gdlpurple!70!black, anchor=north, align=center] at (1.4, 0.2)
    {$m_v=\bigoplus_{u\in\Nbr(v)} m_{u\to v}$};
\end{scope}
\draw[->, line width=1.6pt, gdlgray!70] (8.2, 2.0) -- (9.0, 2.0);
\begin{scope}[shift={(9.2,0)}]
  \node[centernode, minimum size=26pt] (vu) at (1.4, 2.0) {$v$};
  \node[font=\footnotesize, text=gdlgreen!70!black, anchor=north, align=center] at (1.4, 0.6)
    {$h_v'=U_\ell(h_v,m_v)$};
\end{scope}
\end{tikzpicture}
\caption[Message passing on a graph]{The three phases of a message passing layer
for a single node $v$. Each neighbour $u$ emits a message $m_{u\to v}$; the
messages are combined by a permutation-invariant aggregator $\bigoplus$; the node
state is refreshed by the update function $U_\ell$.}
\label{fig:graphs:message-passing}
\end{figure}

\paragraph{Convolutional, attentional, and general flavours.}
Within this framework, the choice of message function defines a spectrum of
expressivity and cost~\citep{bronstein2021geometric}. In the
\emph{convolutional} flavour the message is a fixed, structure-determined
weighting of the neighbour, $M_\ell(h_v, h_u) = c_{vu}\, \psi(h_u)$ with scalar
coefficients $c_{vu}$ depending only on the graph (e.g.\ degrees); this is the
form taken by the graph convolutional network of \citet{kipf2017semi}. In the
\emph{attentional} flavour the coefficient is learned from the endpoints,
$c_{vu} = a(h_v, h_u)$, as in the graph attention network of
\citet{velickovic2017graph} (\Cref{sec:graphs:gat}). In the fully \emph{general}
flavour the message is an arbitrary learnable function of both endpoints, as in
\Cref{def:graphs:mpnn}. Cost and flexibility increase along this spectrum.

\subsection{Convolution as message passing on a grid}
\label{sec:graphs:cnn-mpnn}

The power of message passing is that it subsumes the architectures of the
preceding chapters. \Cref{ch:grids} treated the convolutional case in detail; we
make the identification precise here.

\begin{theorem}[A CNN is a special case of a GNN]
\label{thm:graphs:cnn-gnn}
A 2D convolutional network is exactly equivalent to a graph neural network whose
graph is the regular grid $G = (\Z^2 \cap \Omega, E)$, whose vertices are the
pixels $(i,j)$, whose edges encode $4$- or $8$-neighbour connectivity, and whose
weights are shared according to relative position (the kernel). Concretely, with
the message $M_\ell(h_v, h_u) = W_{v-u}\, h_u$, sum aggregation
$m_v = \sum_{u \in \Nbr(v)} M_\ell$ and update
$U_\ell(h_v, m_v) = \sigma(m_v + b)$, the node states of the GNN equal the
activations of the CNN at every layer.
\end{theorem}

\begin{proof}
By induction on depth. \emph{Base case:} $h_v^{(0)} = x_{i,j} = y_{i,j}^{(0)}$ by
construction, where $v = (i,j)$. \emph{Inductive step:} assume
$h_v^{(\ell)} = y_{i,j}^{(\ell)}$. The CNN output at $(i,j)$ in layer $\ell+1$ is
\[
  y_{i,j}^{(\ell+1)}
    = \sigma\!\Big(\sum_{(u,v)\in\mathcal{K}} K_{u,v}\, y_{i+u,j+v}^{(\ell)} + b\Big)
    = \sigma\!\Big(\sum_{(i',j')\in\Nbr(i,j)} K_{i'-i,j'-j}\, y_{i',j'}^{(\ell)} + b\Big),
\]
substituting $(i',j') = (i+u, j+v)$. The GNN update gives the same expression
with $W_{v-u} = K_{i'-i,j'-j}$, equal to the CNN by the inductive hypothesis, so
$h_v^{(\ell+1)} = y_{i,j}^{(\ell+1)}$.
\end{proof}

The crucial difference between a convolution and a general graph layer is the
\emph{weight sharing} $W_{v-u}$: the message weight depends only on the relative
position $v-u$, which is meaningful solely on a domain with translation
structure. On a grid this enforces translation equivariance and collapses the
parameter count from $O(n^2)$ to $O(\abs{\mathcal{K}})$. On an arbitrary graph
there is no notion of relative position, so the message weight must instead be a
function of the node features (attention) or of structural invariants such as
degree (convolution).

\begin{example}[Numerical verification: the $3\times3$ Sobel kernel]
\label{ex:graphs:sobel}
Consider the horizontal Sobel kernel and a $3 \times 3$ image,
\[
  K = \begin{bmatrix} -1 & 0 & 1 \\ -2 & 0 & 2 \\ -1 & 0 & 1 \end{bmatrix},
  \qquad
  I = \begin{bmatrix} 0 & 1 & 2 \\ 0 & 1 & 2 \\ 0 & 1 & 2 \end{bmatrix}.
\]
\textbf{CNN method.} The convolution at the central pixel $(1,1)$ gives
\[
  y_{1,1} = (-1)(0) + (1)(2) + (-2)(0) + (2)(2) + (-1)(0) + (1)(2) = 8 .
\]
\textbf{GNN method.} Building the grid graph with $v_{1,1}$ connected to its
eight neighbours, the aggregated message is
\[
  m_{1,1} = \sum_{u} K_{v-u}\, h_u
    = (-1)(0) + (1)(2) + (-2)(0) + (2)(2) + (-1)(0) + (1)(2) = 8 .
\]
Both methods produce the identical result $y_{1,1} = 8$, verifying the
equivalence of \Cref{thm:graphs:cnn-gnn}.
\end{example}

\begin{figure}[h]
\centering
\begin{tikzpicture}[>=Stealth, font=\small,
    pix/.style={circle, draw=gdlblue!60, fill=gdlblue!15, minimum size=16pt,
        inner sep=0pt, font=\scriptsize},
    cpix/.style={circle, draw=gdlred!70, fill=gdlred!20, minimum size=18pt,
        inner sep=0pt, font=\scriptsize\bfseries}]
\node[font=\bfseries, text=gdlorange!70!black] at (1,3.5) {Convolution};
\node[font=\bfseries, text=gdlgreen!70!black] at (8,3.5) {Message passing};
% left: 3x3 grid + kernel
\foreach \i in {0,1,2} \foreach \j in {0,1,2}{
  \ifnum\i=1 \ifnum\j=1 \node[cpix] (g\i\j) at (\j,2-\i){}; \else
    \node[pix] (g\i\j) at (\j,2-\i){}; \fi
  \else \node[pix] (g\i\j) at (\j,2-\i){}; \fi}
\foreach \i in {0,1,2} \draw[gdlgray!50] (g\i0)--(g\i1)--(g\i2);
\foreach \j in {0,1,2} \draw[gdlgray!50] (g0\j)--(g1\j)--(g2\j);
\draw[gdlgray!50] (g00)--(g11)--(g22) (g02)--(g11)--(g20);
\node[font=\scriptsize] at (1,-0.7) {$3\times3$ neighbourhood};
\draw[->, line width=1.4pt, gdlgray!70] (3,1) -- (4.6,1)
  node[midway,above,font=\scriptsize]{$=$};
% right: central node with messages
\node[cpix, minimum size=22pt] (cc) at (8,1) {$v$};
\foreach \a in {45,90,135,180,225,270,315,0}{
  \node[pix] (m\a) at ($(8,1)+(\a:1.4)$){};
  \draw[->, gdlgreen!70, line width=1pt, shorten >=3pt, shorten <=3pt] (m\a)--(cc);}
\node[font=\scriptsize, anchor=north, align=center] at (8,-0.5)
  {$m_v=\sum_{u\in\Nbr(v)} W_{v-u} h_u$};
\end{tikzpicture}
\caption[Convolution as message passing]{Convolution equals message passing on a
regular grid: the three phases (message, aggregation, update) coincide exactly
with the convolutional operation when the graph is a regular grid.}
\label{fig:graphs:conv-mp}
\end{figure}

For a $H \times W$ image with $C$ channels and a $k \times k$ kernel, a standard
CNN costs $O(HW k^2 C)$ while a grid GNN costs $O(\abs{E}\, C) = O(4HW\, C)$ for
$4$-connectivity; the two coincide when $k$ is small. The GNN viewpoint reveals
that complexity depends on $\abs{E}$, not on kernel size, giving a unified lens
for analysing efficiency across topologies. Analogously, self-attention is
message passing on the complete graph $K_n$ with content-dependent coefficients
$\alpha_{ij} = \operatorname{softmax}_j\big((Q h_i)^\top (K h_j)/\sqrt{d_k}\big)$
and identity update --- the attentional flavour applied to an all-to-all
topology. Thus a fully connected layer, a convolution, a transformer, and a graph
network differ only in the graph topology and in how the message weight is
determined; they are instances of one geometric construction.

\subsection{From spectral CNNs to spectral GNNs}
\label{sec:graphs:cheb}

The spectral viewpoint provides the formal bridge between traditional CNNs and
graph networks, showing both to be special cases of the same principle.

\begin{theorem}[Fundamental CNN--GNN connection]
\label{thm:graphs:cnn-spectral}
Every CNN can be represented exactly as a spectral GNN over a regular grid: the
graph $G = (V, E)$ represents the pixel structure, connectivity $E$ is given by
spatial neighbourhood, and the spectral filters correspond to convolutional
kernels via the discrete Fourier transform.
\end{theorem}

\begin{proof}[Proof sketch]
For an $n \times n$ grid the Laplacian has entries $4$ on the diagonal, $-1$ for
$\abs{i-i'}+\abs{j-j'} = 1$, and $0$ otherwise. Its eigenvectors are exactly the
two-dimensional discrete Fourier functions
$\phi_{k,\ell}(i,j) = \frac1n e^{2\pi\imath(ki+\ell j)/n}$. Hence spatial
convolution $f \star g$ corresponds exactly to spectral multiplication
$\hat f \odot \hat g$ in the graph frequency domain.
\end{proof}

All three drawbacks of naive spectral filtering are removed by restricting the
transfer function to a \emph{polynomial} of the Laplacian. If
$g_\theta(\lambda) = \sum_{k=0}^K \theta_k \lambda^k$, then
$g_\theta \star f = (\sum_k \theta_k \Lap^k) f$ and the eigenvectors cancel out:
no eigendecomposition is needed. Moreover $\Lap^k$ connects only nodes within
graph distance $k$, so a degree-$K$ polynomial filter is exactly $K$-localized.
This is the precise sense in which the spectral and spatial views coincide.

\begin{definition}[ChebNet: Chebyshev convolution]
\label{def:graphs:chebnet}
A spectral filter approximated by Chebyshev polynomials of order $K$ is
$g_\theta(\Lambda) = \sum_{k=0}^K \theta_k T_k(\tilde\Lambda)$, where $T_k$ are
the Chebyshev polynomials of the first kind and
$\tilde\Lambda = 2\Lambda/\lambda_{\max} - I$ is the rescaled Laplacian
\citep{defferrard2017convolutionalneuralnetworksgraphs}.
\end{definition}

\begin{theorem}[Localization of Chebyshev filters]
\label{thm:graphs:cheb-local}
A Chebyshev filter of order $K$ is strictly $K$-localized: the response at node
$i$ depends only on nodes at graph distance at most $K$ from $i$.
\end{theorem}

\begin{proof}
The Chebyshev recurrence is
$T_k(\tilde\Lap) = 2\tilde\Lap\, T_{k-1}(\tilde\Lap) - T_{k-2}(\tilde\Lap)$. Since
$\tilde\Lap_{ij} \neq 0$ only for adjacent $i, j$, multiplication by $\tilde\Lap$
extends the support of a signal by exactly one hop. Thus $T_1(\tilde\Lap) f$
depends on neighbours at distance $\le 1$, and by induction $T_k(\tilde\Lap) f$
at node $i$ depends only on nodes at distance $\le k$. As
$g_\theta(\tilde\Lap) = \sum_{k=0}^K \theta_k T_k(\tilde\Lap)$, the response
depends on nodes at distance $\le K$.
\end{proof}

The recurrence lets one compute $g_\theta(\tilde\Lap) f$ by repeated sparse
matrix--vector products, at total cost $O(K\abs{E})$ and with only $K+1$
parameters, none depending on the eigenbasis. The graph convolutional network of
\citet{kipf2017semi} is the first-order case $K = 1$ with a single shared
parameter and a renormalization of the adjacency matrix, giving
$H^{(\ell+1)} = \sigma\!\big(\hat D^{-1/2} \hat A\, \hat D^{-1/2} H^{(\ell)}
W^{(\ell)}\big)$ with $\hat A = A + I$. This is plainly a message passing layer
(\Cref{def:graphs:mpnn}) with degree-normalized aggregation, closing the circle:
the spectral construction, made local and parameter-efficient, \emph{is} spatial
message passing.

\paragraph{A hierarchy of generalization.}
The equivalences above place neural architectures in a single hierarchy of
increasing generality:
\[
  \text{Perceptron} \subset \text{MLP} \subset \text{CNN (grid)}
    \subset \text{spectral CNN} \subset \text{spectral GNN}
    \subset \text{spatial GNN} \subset \text{general MPNN},
\]
each inclusion relaxing one architectural restriction (the spatial locality of
the kernel, then the regular grid, then the fixed graph), at the cost of greater
expressive power, more parameters, and increased overfitting risk. CNNs and GNNs
share three architectural principles --- locality (aggregation from local
neighbourhoods), weight sharing (parameters independent of position), and
hierarchical composition (multi-scale representations via pooling or graph
coarsening) --- which is why techniques such as pooling, dropout, batch
normalization and skip connections transfer freely between them.

\section{Graph attention networks}
\label{sec:graphs:gat}

The attentional flavour of message passing deserves separate treatment. Graph
attention networks (GAT) sit at the intersection of spectral graph theory and
adaptive attention: where spectral methods require a costly eigendecomposition,
GATs achieve comparable expressive power through learnable attention coefficients
that implicitly capture graph structure and long-range dependencies. Attention is
message passing in which the aggregation weights are computed from the node
features themselves --- and, applied to the complete graph, it recovers the
self-attention of \Cref{ch:grids}.

\begin{definition}[GAT attention mechanism]
\label{def:graphs:gat}
For a node $i$ with neighbourhood $\Nbr(i)$, the GAT
mechanism~\citep{velickovic2017graph} computes:
\begin{align}
  z_i &= W h_i, \qquad z_j = W h_j
    && \text{(shared linear projection),} \\
  e_{ij} &= a^\top [z_i \,\|\, z_j]
    && \text{(attention coefficients),} \\
  \alpha_{ij} &= \frac{\exp\!\big(\operatorname{LeakyReLU}(e_{ij})\big)}
       {\sum_{k \in \Nbr(i) \cup \{i\}} \exp\!\big(\operatorname{LeakyReLU}(e_{ik})\big)}
    && \text{(masked softmax),} \\
  h_i' &= \sigma\!\Big(\sum_{j \in \Nbr(i) \cup \{i\}} \alpha_{ij}\, z_j\Big)
    && \text{(weighted aggregation),}
\end{align}
where $W \in \R^{d' \times d}$ is a shared weight matrix, $a \in \R^{2d'}$ a
learnable attention vector, and $\|$ denotes concatenation.
\end{definition}

\begin{figure}[h]
\centering
\begin{tikzpicture}[>=Stealth, font=\small,
    targetnode/.style={circle, draw=gdlred!70, fill=gdlred!20, minimum size=26pt,
        inner sep=0pt, font=\small\bfseries, line width=1.4pt},
    neighnode/.style={circle, draw=gdlblue!60, fill=gdlblue!20, minimum size=20pt,
        inner sep=0pt, font=\small},
    nonneigh/.style={circle, draw=gdlgray!40, fill=gdlgray!15, minimum size=18pt,
        inner sep=0pt, font=\small, text=gdlgray!70!black}]
\node[targetnode] (vi) at (0,0) {$v_i$};
\node[neighnode] (v1) at (-2.2,1.6) {$v_1$};
\node[neighnode] (v2) at (0,2.4) {$v_2$};
\node[neighnode] (v3) at (2.2,1.3) {$v_3$};
\node[neighnode] (v4) at (-2.0,-1.3) {$v_4$};
\node[nonneigh] (v5) at (2.4,-1.5) {$v_5$};
\node[nonneigh] (v6) at (-3.3,0.2) {$v_6$};
\draw[gdlgray!40, thin] (v1)--(v6) (v4)--(v6) (v3)--(v5) (v1)--(v2)
  (vi)--(v1) (vi)--(v2) (vi)--(v3) (vi)--(v4) (v5)--(v3);
\draw[-{Stealth[length=6pt]}, line width=3pt, gdlorange, opacity=0.9,
  shorten >=7pt, shorten <=6pt] (v2)--(vi);
\draw[-{Stealth[length=5pt]}, line width=2pt, gdlorange, opacity=0.75,
  shorten >=7pt, shorten <=6pt] (v1)--(vi);
\draw[-{Stealth[length=5pt]}, line width=1.4pt, gdlorange, opacity=0.6,
  shorten >=7pt, shorten <=6pt] (v3)--(vi);
\draw[-{Stealth[length=4pt]}, line width=0.8pt, gdlorange, opacity=0.45,
  shorten >=7pt, shorten <=6pt] (v4)--(vi);
\node[font=\scriptsize, text=gdlorange!70!black] at (-0.7,1.4) {$\alpha_{2i}=0.42$};
\node[font=\scriptsize, text=gdlorange!70!black] at (-1.7,0.5) {$\alpha_{1i}=0.28$};
\node[font=\scriptsize, text=gdlorange!70!black] at (1.5,0.5) {$\alpha_{3i}=0.19$};
\node[font=\scriptsize, text=gdlorange!70!black] at (-1.5,-0.7) {$\alpha_{4i}=0.11$};
\node[rounded corners=4pt, draw=gdlorange!60, fill=gdlorange!8, line width=1pt,
  inner sep=7pt, font=\footnotesize, text=gdlorange!70!black, anchor=north west,
  align=left] at (3.4,1.0)
  {\textbf{Attention coefficients:}\\[3pt]
   $e_{ij}=\operatorname{LeakyReLU}(a^\top[Wh_i\,\|\,Wh_j])$\\[4pt]
   $\alpha_{ij}=\operatorname{softmax}_j(e_{ij})$\\[4pt]
   $h_i'=\sigma\big(\sum_{j\in\Nbr(i)}\alpha_{ij}Wh_j\big)$};
\node[font=\scriptsize\bfseries, text=gdlgray!70!black, anchor=west]
  at (-3.3,-2.4) {Arrow width $\propto$ attention weight};
\end{tikzpicture}
\caption[GAT attention mechanism]{The graph attention mechanism. The central
node (red) computes attention coefficients $\alpha_{ij}$ with each neighbour
(blue). Arrow thickness reflects the magnitude of the attention weights: the
network learns to assign different importance to each neighbour based on node
features.}
\label{fig:graphs:gat}
\end{figure}

\begin{proposition}[Expressive power of GAT]
\label{prop:graphs:gat}
Standard graph attention networks are bounded in expressive power by the $1$-\WL{}
test (\Cref{thm:graphs:wl-mpnn}). In particular: GAT cannot distinguish pairs of
non-isomorphic graphs that $1$-\WL{} cannot; unlike a graph convolutional network
with fixed weights, GAT can learn \emph{anisotropic}, adaptive aggregation
functions through its attention coefficients; and GAT is permutation equivariant
by construction.
\end{proposition}

\begin{proof}
GAT is an instance of the MPNN framework with attention-weighted aggregation over
the $1$-hop neighbourhood, so by \Cref{thm:graphs:wl-mpnn} it is bounded by
$1$-\WL. Permutation equivariance follows because the coefficients $\alpha_{ij}$
depend only on $(h_i, h_j)$ and the aggregation
$\sum_{j \in \Nbr(i)} \alpha_{ij} z_j$ is symmetric in the order of the
neighbours.
\end{proof}

To improve stability and expressive capacity, GATs employ \emph{multi-head
attention}: $h$ independent attention mechanisms, each with its own parameters
$W^{(k)}, a^{(k)}$, run in parallel and are combined by averaging,
$h_i' = \sigma\big(\tfrac1h \sum_{k=1}^h \sum_{j \in \Nbr(i)} \alpha_{ij}^{(k)}
W^{(k)} h_j\big)$, or by concatenation. Multiple heads reduce gradient variance,
let different heads specialize in different relational patterns, and enlarge the
representable function class.

\paragraph{Spectral positional encoding.}
A fundamental limitation of GAT (and of message passing generally) is its
inability to distinguish structurally equivalent nodes. \emph{Spectral positional
encoding} resolves this by attaching to each node a signature drawn from the
Laplacian eigenvectors. With $L = U\Lambda U^\top$, the encoding of node $i$ is
$\operatorname{PE}_i = [u_{i,1}, \dots, u_{i,k}] \in \R^k$, where $u_{i,j}$ is the
$i$-th entry of the $j$-th smallest nontrivial eigenvector
(\Cref{def:graphs:gft}). One may further define \emph{spectral distances} for
attention between distant nodes --- the diffusion distance
$d_D^2(i,j) = \sum_{k\ge2} e^{-2t\lambda_k}(u_{i,k} - u_{j,k})^2$, the Green
function $G(i,j) = \sum_{k\ge2} u_{i,k} u_{j,k}/\lambda_k$, and the biharmonic
distance $d_B^2(i,j) = \sum_{k\ge2}(u_{i,k} - u_{j,k})^2/\lambda_k^2$. Spectral
graph transformers~\citep{kreuzer2021spectral} combine such learnable positional
encodings with full attention, giving every pair of nodes a direct connection and
thereby resolving the oversquashing pathology (\Cref{sec:graphs:oversquashing})
at $O(1)$ information distance.

\section{Expressivity and the Weisfeiler--Leman hierarchy}
\label{sec:graphs:wl}

Having a unifying mechanism invites a sharp question: how powerful is it? Can a
message passing network tell any two non-isomorphic graphs apart? The answer is
no, and the precise limit is a classical combinatorial algorithm from graph
isomorphism testing, originally proposed in 1968~\citep{weisfeiler1968reduction}.

\subsection{The Weisfeiler--Leman test}
\label{sec:graphs:wl-test}

\begin{definition}[1-Weisfeiler--Leman test]
\label{def:graphs:wl}
Given a graph $G = (V, E)$, the \emph{$1$-Weisfeiler--Leman} ($1$-\WL) algorithm
assigns colours to nodes and refines them iteratively. Initialize $c^{(0)}(v)$
identically for all nodes (or from node labels, if present). At each step,
\[
  c^{(t+1)}(v) = \HASH\!\Big(c^{(t)}(v),\,
      \{\!\{\, c^{(t)}(u) : u \in \Nbr(v) \,\}\!\}\Big),
\]
where $\{\!\{\cdot\}\!\}$ is a multiset and $\HASH$ is an injective map sending
each distinct (colour, neighbour-multiset) pair to a fresh colour. The refinement
stabilizes in at most $n$ steps; two graphs are \emph{$\WL$-equivalent}
($G_1 \equiv_{\WL} G_2$) when the final colour histograms coincide.
\end{definition}

The structural resemblance between the refinement step and a message passing
layer is unmistakable: a node updates its colour from its own colour and the
multiset of its neighbours' colours, exactly as a node updates its features from
its own features and the aggregate of its neighbours' messages.
\Cref{fig:graphs:wl} illustrates the refinement.

\begin{figure}[h]
\centering
\begin{tikzpicture}[>=Stealth, font=\small,
    wlnode/.style={circle, draw, thick, minimum size=6mm, inner sep=0pt}]
\foreach \dx in {0,5.0,10.0}{
  \begin{scope}[shift={(\dx,0)}]
    \coordinate (n1) at (0,1.6);\coordinate (n2) at (1.8,1.6);
    \coordinate (n3) at (1.8,0);\coordinate (n4) at (0,0);
    \coordinate (n5) at (0.9,-1.1);\coordinate (n6) at (0.9,0.8);
    \draw[thick,gdlgray!50] (n1)--(n2) (n2)--(n3) (n3)--(n4) (n4)--(n1)
       (n3)--(n5) (n4)--(n5) (n1)--(n6) (n2)--(n6) (n3)--(n6) (n4)--(n6);
  \end{scope}
}
\begin{scope}
  \node at (0.9,2.6) {\bfseries (a) $c^{(0)}$};
  \foreach \p/\x/\y in {1/0/1.6,2/1.8/1.6,3/1.8/0,4/0/0,5/0.9/-1.1,6/0.9/0.8}
    \node[wlnode, fill=gdlgray!40] at (\x,\y) {};
\end{scope}
\draw[->, very thick, gdlgray!50] (2.7,0.6) -- node[above,font=\scriptsize]{HASH} (4.3,0.6);
\begin{scope}[shift={(5.0,0)}]
  \node at (0.9,2.6) {\bfseries (b) $c^{(1)}$};
  \foreach \p/\x/\y/\c in {1/0/1.6/{gdlblue!50},2/1.8/1.6/{gdlblue!50},3/1.8/0/{gdlred!50},4/0/0/{gdlred!50},5/0.9/-1.1/{gdlgreen!50},6/0.9/0.8/{gdlred!50}}
    \node[wlnode, fill=\c] at (\x,\y) {};
\end{scope}
\draw[->, very thick, gdlgray!50] (7.7,0.6) -- node[above,font=\scriptsize]{HASH} (9.3,0.6);
\begin{scope}[shift={(10.0,0)}]
  \node at (0.9,2.6) {\bfseries (c) $c^{(2)}$};
  \foreach \p/\x/\y/\c in {1/0/1.6/{gdlblue!50},2/1.8/1.6/{gdlblue!50},3/1.8/0/{gdlyellow!60},4/0/0/{gdlyellow!60},5/0.9/-1.1/{gdlgreen!50},6/0.9/0.8/{gdlred!50}}
    \node[wlnode, fill=\c] at (\x,\y) {};
\end{scope}
\end{tikzpicture}
\caption[Weisfeiler--Leman colour refinement]{The $1$-\WL{} refinement on a
six-node graph. (a)~Uniform initial colouring. (b)~After one step nodes are
distinguished by degree. (c)~A further step separates the two degree-$4$ nodes
$v_3, v_4$ from $v_6$, because their neighbour-colour multisets differ.}
\label{fig:graphs:wl}
\end{figure}

This resemblance is made into a theorem by \citet{xu2019how} and
\citet{morris2019weisfeiler}.

\begin{theorem}[Expressive power of MPNNs equals 1-\WL]
\label{thm:graphs:wl-mpnn}
Message passing neural networks are bounded above by $1$-\WL: if two graphs are
$\WL$-equivalent, every MPNN assigns them identical representations. Moreover the
bound is tight: there exist MPNNs (the Graph Isomorphism Network of
\citet{xu2019how}) that distinguish exactly the graphs $1$-\WL{} distinguishes,
achieved by using injective aggregation and update functions.
\end{theorem}

\begin{proof}[Proof sketch]
For the upper bound, an MPNN layer computes node states from the node's own state
and the multiset of neighbour states --- the same information $1$-\WL{} uses; by
induction on layers, $\WL$-equivalent graphs map to equal representations. For
tightness, one builds a message and update realizing an injective hash: sum
aggregation over a countable feature space followed by a multilayer perceptron is
injective on multisets by the universal approximation theorem for set
functions~\citep{zaheer2017deepsets}, so the MPNN reproduces each refinement step
exactly.
\end{proof}

\begin{example}[Two graphs $1$-\WL{} cannot separate]
\label{ex:graphs:wl-fail}
Let $G_1$ be the six-cycle $C_6$ and $G_2$ the disjoint union of two triangles
$K_3 \sqcup K_3$. Both have six nodes, six edges, and every node of degree $2$.
Under $1$-\WL{} all nodes start with one colour, see the multiset
$\{\!\{1,1\}\!\}$, and stabilize to a single colour; the histograms are identical
($\mathbf{h}(G_1) = \mathbf{h}(G_2) = (6)$), so $G_1 \equiv_{\WL} G_2$. Yet the
graphs are not isomorphic: $C_6$ is connected while $K_3 \sqcup K_3$ has two
components. A standard MPNN therefore cannot tell a hexagonal ring from a pair of
triangles, an instructive failure when global connectivity matters --- for
example, distinguishing certain structural isomers in chemistry.
\end{example}

\subsection{Oversmoothing}
\label{sec:graphs:oversmoothing}

\Cref{thm:graphs:wl-mpnn} bounds what a network can distinguish but says nothing
about what happens as we stack many layers. Empirically, deep MPNNs degrade: node
representations converge to a common value and lose all discriminative power.
This is \emph{oversmoothing}~\citep{li2018deeper}, and it has a clean spectral
explanation through the Dirichlet energy.

\begin{definition}[Dirichlet energy]
\label{def:graphs:dirichlet}
For node features $H \in \R^{n \times d}$ the \emph{Dirichlet energy} is
\[
  E(H) = \tr(H^\top L\, H)
       = \sum_{c=1}^d \sum_{(i,j) \in E} W_{ij}\,(h_{ic} - h_{jc})^2 ,
\]
which vanishes iff $H$ is constant on each connected component.
\end{definition}

\begin{proposition}[Linear MPNN as discrete diffusion]
\label{prop:graphs:diffusion}
A linear, mean-aggregating message passing step
$h_v^{(\ell+1)} = (1-\alpha) h_v^{(\ell)} + \alpha \sum_{u \in \Nbr(v)}
h_u^{(\ell)}/\deg(v)$ with $\alpha \in (0,1]$ is, in matrix form,
$H^{(\ell+1)} = (I - \alpha \tilde L)\, H^{(\ell)}$, with the random-walk
Laplacian $\tilde L = D^{-1} L = I - D^{-1} A$, which shares its spectrum with the
symmetric normalized Laplacian.
\end{proposition}

\begin{proof}
For each node,
$h_v^{(\ell+1)} = h_v^{(\ell)} - \alpha\big(h_v^{(\ell)}
- \sum_{u \sim v} h_u^{(\ell)}/\deg(v)\big)
= h_v^{(\ell)} - \alpha (D^{-1} L\, H^{(\ell)})_v
= (I - \alpha\tilde L\, H^{(\ell)})_v$, using
$(D^{-1} A H)_v = \sum_{u \sim v} h_u/\deg(v)$.
\end{proof}

This is the explicit Euler discretization of the heat equation
$\partial_t H = -\alpha \tilde L H$ on the graph, and iterating it smooths the
signal: $H^{(\ell)} = (I - \alpha\tilde L)^\ell H^{(0)}$.

\begin{theorem}[Exponential oversmoothing]
\label{thm:graphs:oversmoothing}
Let $G$ be connected with normalized-Laplacian eigenvalues
$0 = \lambda_1 < \lambda_2 \le \cdots \le \lambda_n$ and let
$0 < \alpha < 1/\lambda_n$. Then after $\ell$ linear diffusion layers
\[
  E\big(H^{(\ell)}\big) \le (1 - \alpha \lambda_2)^{2\ell}\, E\big(H^{(0)}\big),
\]
so $E(H^{(\ell)}) \to 0$ exponentially: all node representations collapse to a
common vector.
\end{theorem}

\begin{proof}
Let $\{u_k\}$ be the orthonormal eigenbasis of
$L_{\mathrm{sym}} = D^{-1/2} L\, D^{-1/2}$ with eigenvalues $\lambda_k$, and set
$f_c = D^{1/2} h_c$ for each column of $H$, so that
$h_c^\top L\, h_c = f_c^\top L_{\mathrm{sym}} f_c$. The diffusion operator
satisfies $D^{1/2}(I - \alpha\tilde L)D^{-1/2} = I - \alpha L_{\mathrm{sym}}$,
scaling the coefficient of $u_k$ by $(1 - \alpha\lambda_k)$. Hence
\[
  E(H^{(\ell)}) = \sum_c \sum_{k \ge 2}
    \lambda_k (1-\alpha\lambda_k)^{2\ell} \abs{\hat f_{c,k}}^2
\]
(the $k=1$ term vanishes since $\lambda_1 = 0$). Since $\alpha < 1/\lambda_n$,
$\abs{1-\alpha\lambda_k} \le 1 - \alpha\lambda_2$ for all $k \ge 2$, which yields
the bound; as $0 < \alpha\lambda_2 < 1$, the factor decays exponentially.
\end{proof}

The rate is governed by the \emph{spectral gap} $\lambda_2$: well-connected
graphs (large $\lambda_2$, such as expanders) homogenize fastest, while poorly
connected graphs ($\lambda_2 \approx 0$) preserve diversity longer. This is the
discrete shadow of heat diffusion on a Riemannian manifold, where the spectral
gap of the Laplace--Beltrami operator controls the mixing rate. Residual
connections, layer normalization, and sheaf-based diffusion~\citep{bodnar2022neural}
--- which replaces scalar diffusion by diffusion in sheaves, letting each edge
transport information in a different vector space --- are among the remedies.

\subsection{Oversquashing and higher-order tests}
\label{sec:graphs:oversquashing}

If oversmoothing is the loss of high-frequency signal as depth grows, the
complementary failure is \emph{oversquashing}: the inability to propagate
information across long distances, because the messages from an exponentially
growing $\ell$-hop neighbourhood must be compressed into a fixed-size vector.

\begin{definition}[Jacobian sensitivity]
\label{def:graphs:sensitivity}
The \emph{sensitivity} of node $v$ to node $u$ after $\ell$ layers is
$S_{vu}^{(\ell)} = \norm{\partial h_v^{(\ell)} / \partial h_u^{(0)}}$, the
operator norm of the Jacobian $\partial h_v^{(\ell)} / \partial h_u^{(0)} \in
\R^{d \times d}$. If $S_{vu}^{(\ell)} \approx 0$, node $v$ is essentially
insensitive to the input of node $u$.
\end{definition}

\begin{proposition}[Exponential decay of sensitivity]
\label{prop:graphs:oversquashing}
For an MPNN with message and update functions Lipschitz with constants
$c_M, c_U$ and $\hat A = D^{-1} A$ the random-walk transition matrix,
\[
  S_{vu}^{(\ell)} \le (c_U c_M)^\ell\, [\hat A^\ell]_{vu}.
\]
In particular, if the graph distance $d_G(v,u) > \ell$, then
$[\hat A^\ell]_{vu} = 0$ and the sensitivity is exactly zero.
\end{proposition}

\begin{proof}[Proof sketch]
By induction on $\ell$. For $\ell = 1$ and $v \neq u$, the chain rule gives
$\partial h_v^{(1)} / \partial h_u^{(0)}
= (\partial U / \partial m_v)(\partial m_v / \partial h_u^{(0)})$, nonzero only if
$u \in \Nbr(v)$, where the Lipschitz constant contributes $c_M$ and the degree
normalization the factor $1/\deg(v) = [\hat A]_{vu}$; hence
$S_{vu}^{(1)} \le c_U c_M [\hat A]_{vu}$. For the inductive step, the multivariate
chain rule over $w \in \Nbr(v) \cup \{v\}$ together with the inductive hypothesis
produces the matrix product $\hat A^{\ell+1}$, recording all length-$(\ell+1)$
paths from $u$ to $v$.
\end{proof}

The bound shows that the graph topology, encoded in $\hat A^\ell$, is the
determining factor. \citet{topping2022understanding} tie these bottlenecks to
negative discrete (Ollivier) Ricci curvature on edges: edges with negative Ricci
curvature are bottlenecks where $[\hat A^\ell]_{vu}$ decays sharply. This is the
discrete version of the Jacobi equation in Riemannian geometry (\Cref{ch:geodesics}):
negative sectional curvature makes geodesics diverge, reducing the volume of
information that can propagate. They propose graph \emph{rewiring} via Ricci flow,
adding edges in negatively curved regions to alleviate the bottlenecks --- a clean
instance of differential geometry prescribing a cure for a learning pathology.

Finally, the $1$-\WL{} barrier can be raised by climbing a hierarchy of more
powerful tests.

\begin{definition}[$k$-Weisfeiler--Leman test]
\label{def:graphs:k-wl}
The \emph{$k$-\WL} test colours $k$-tuples of nodes
$\bar v = (v_1, \dots, v_k) \in V^k$, refining each tuple's colour from the
colours of all tuples obtained by replacing one of its entries by an arbitrary
node~\citep{morris2019weisfeiler}:
\[
  c^{(t+1)}(\bar v) = \HASH\!\Big(c^{(t)}(\bar v),\,
    \{\!\{\, (c^{(t)}(\bar v[w/1]), \dots, c^{(t)}(\bar v[w/k])) : w \in V \,\}\!\}\Big),
\]
where $\bar v[w/i]$ replaces the $i$-th component of $\bar v$ by $w$.
\end{definition}

\begin{proposition}[Strict $k$-\WL{} hierarchy]
\label{prop:graphs:k-wl}
The tests form a strictly increasing hierarchy in distinguishing power,
\[
  1\text{-}\WL \subsetneq 2\text{-}\WL \subsetneq 3\text{-}\WL \subsetneq \cdots,
\]
each level distinguishing graphs the previous cannot, at a cost of $O(n^k)$
memory per iteration ($O(n^k \log n)$ time).
\end{proposition}

The exponential cost of $k$-\WL{} motivates cheaper ways past the $1$-\WL{}
ceiling: structural and positional encodings (for instance Laplacian eigenvectors
from \Cref{def:graphs:gft}) that break node symmetries, spectral graph
transformers~\citep{kreuzer2021spectral}, and message passing on higher-order
topological domains (simplicial and cell complexes, \Cref{sec:graphs:tdl}),
pursued further in \Cref{ch:gauges,ch:applications}.

\medskip
The three limitations --- oversmoothing, oversquashing, and the $\WL$ barrier ---
are facets of one fact: message passing is fundamentally \emph{local}, moving
information one hop per layer. Too many layers destroy the signal; too few prevent
long-range communication; and purely local computation cannot see certain global
patterns. The geometric tools developed here --- spectral analysis, discrete
curvature, higher-order topology --- do not merely diagnose these limits but
prescribe concrete remedies, exemplifying the thesis of this book that geometry is
the organizing principle of effective architectures.
